\chapter{Complex Manifolds}\label{complex-manifolds}
\chapterstatus{complete}

\section{Introduction}\label{complex-manifolds:section-introduction}

A complex manifold is a manifold with holomorphic transition functions. The rigidity of holomorphic functions makes these spaces
very different from smooth manifolds: compact complex manifolds have only constant holomorphic functions, and their geometry is
recorded instead by the cohomology of holomorphic bundles and by the decomposition of differential forms into types, which is the
subject of Hodge theory. This chapter sets up complex manifolds, meromorphic forms on Riemann surfaces and the residue theorem on a compact Riemann surface (the sum of the residues of a
meromorphic $1$-form is zero), the uniformisation of Riemann surfaces with its consequences (the hyperbolic metric, Picard's theorem),
the operators $\partial$ and $\bar\partial$, analytic subvarieties and
blow-ups. References: \cite[Chapter 0]{griffiths-harris}, \cite[Chapter 2]{huybrechts-complex},
\cite[Chapter 3]{voisin-hodge-1}.

\noindent\textbf{Prerequisites.} Chapter~\ref{smooth-manifolds} (Smooth Manifolds), Chapter~\ref{several-complex-variables} (Several
Complex Variables), Chapter~\ref{multilinear-algebra} (Multilinear Algebra) and Chapter~\ref{vector-bundles} (Vector Bundles).

\section{Complex manifolds}\label{complex-manifolds:section-manifolds}

\begin{definition}\label{complex-manifolds:definition-complex-manifold}
A \emph{complex manifold} of dimension $n$ is a topological $2n$-manifold with an atlas of charts $\varphi_i \colon U_i \to \CC^n$ whose transition maps are holomorphic; a map of complex manifolds is
\emph{holomorphic} if it is so in charts, and a \emph{biholomorphism} is a bijective holomorphic map with holomorphic inverse. The sheaf of holomorphic functions is $\mathcal{O}_X$, making $(X, \mathcal{O}_X)$ a
locally ringed space. Every complex manifold is in particular a smooth $2n$-manifold, canonically oriented by the orientation of $\CC^n$ (transition maps have Jacobian determinant $|\det_\CC|^2 > 0$).
\end{definition}

\begin{example}\label{complex-manifolds:example-complex-manifolds}
\begin{enumerate}
\item Open subsets of $\CC^n$; $\CC\PP^n$ with the charts of Example~\ref{smooth-manifolds:example-manifolds}, whose transition maps are rational and nowhere singular.
\item A \emph{complex torus} $X = \CC^n/\Lambda$ for a lattice $\Lambda \cong \ZZ^{2n}$; these are compact, and for $n = 1$ they are the elliptic curves. Two complex tori are rarely biholomorphic even when
they are diffeomorphic, which is the first appearance of moduli (Chapter~\ref{hodge-structures}).
\item Grassmannians $\mathrm{Gr}(k, n)$ of $k$-planes in $\CC^n$, with charts given by graphs of linear maps; they are compact complex manifolds of dimension $k(n-k)$.
\item A smooth analytic hypersurface $\{f = 0\} \subseteq \CC\PP^n$, for $f$ a homogeneous polynomial with $df \neq 0$ on the zero set, is a compact complex manifold of dimension $n - 1$, by the implicit
function theorem (Theorem~\ref{real-analysis:theorem-implicit-function}, valid for holomorphic maps by Remark~\ref{real-analysis:remark-constant-rank}).
\item Not every even-dimensional smooth manifold admits a complex structure: $S^4$ does not (Exercise~\ref{characteristic-classes:exercise-no-almost-complex}), and whether $S^6$ does is open.
\end{enumerate}
\end{example}

\begin{proposition}\label{complex-manifolds:proposition-no-functions}
A holomorphic function on a compact connected complex manifold is constant. Consequently a compact complex manifold has no nonconstant holomorphic map to $\CC^n$, and $\mathcal{O}_X(X) = \CC$.
\end{proposition}

\begin{proof}
$|f|$ attains a maximum at some $p$; in a chart around $p$, the maximum principle
(Proposition~\ref{several-complex-variables:proposition-identity-maximum}) makes $f$ constant near $p$, and the set where $f$ equals $f(p)$ is open and closed by the identity theorem.
\end{proof}

\begin{corollary}\label{complex-manifolds:corollary-compact-submanifolds}
Every compact complex submanifold of $\CC^n$ is a finite set of points. In particular no compact complex manifold of positive dimension, such as $\CC\PP^m$ or a complex torus, embeds holomorphically in $\CC^n$, in contrast with the
Whitney embedding theorem for smooth manifolds (Theorem~\ref{smooth-manifolds:theorem-whitney}).
\end{corollary}

\begin{proof}
The coordinate functions restrict to holomorphic functions on each connected component, which are constant by Proposition~\ref{complex-manifolds:proposition-no-functions}; so each component is a point, and a compact discrete set is
finite.
\end{proof}

\begin{theorem}\label{complex-manifolds:theorem-maps-of-tori}
Let $X = \CC^n/\Lambda$ and $Y = \CC^m/\Lambda'$ be complex tori. Every holomorphic map $f \colon X \to Y$ is induced by an affine map $\CC^n \to \CC^m$, $z \mapsto Az + b$, with $A$ a complex $m \times n$ matrix such that
$A\Lambda \subseteq \Lambda'$. In particular $f$ is a group homomorphism followed by a translation, and $X \cong Y$ as complex manifolds if and only if $n = m$ and $\Lambda' = A\Lambda$ for some $A \in \GL_n(\CC)$.
\end{theorem}

\begin{proof}
The projection $\CC^n \to X$ is a universal covering, as $\CC^n$ is simply connected and $\Lambda$ acts freely and properly (Theorem~\ref{topological-manifolds:theorem-covering-action}). By the lifting criterion
(Proposition~\ref{homotopy:proposition-lifting-criterion}), $f \circ \pi$ lifts to a continuous map $F \colon \CC^n \to \CC^m$, holomorphic because $\pi'$ is a local biholomorphism. For $\lambda \in \Lambda$, $F(z + \lambda) - F(z) \in \Lambda'$ for all
$z$, and by continuity and discreteness of $\Lambda'$ it is a constant $c(\lambda)$. Hence each partial derivative $\partial F_j/\partial z_k$ is $\Lambda$-periodic, so it is a holomorphic function on the compact manifold $X$, hence constant
(Proposition~\ref{complex-manifolds:proposition-no-functions}). So $F(z) = Az + b$, and $A\lambda = c(\lambda) \in \Lambda'$. If $f$ is a biholomorphism, apply this to $f$ and $f^{-1}$: their lifts compose to a lift of the identity, which is
a translation by an element of $\Lambda$, so $A$ is invertible with $A\Lambda = \Lambda'$ and $n = m$; the converse is clear.
\end{proof}

\begin{corollary}\label{complex-manifolds:corollary-elliptic-curves}
Every one-dimensional complex torus is biholomorphic to $E_\tau = \CC/(\ZZ + \ZZ\tau)$ with $\operatorname{Im}\tau > 0$, and $E_\tau \cong E_{\tau'}$ if and only if $\tau' = \frac{a\tau + b}{c\tau + d}$ for some
$\begin{pmatrix}a & b\\ c & d\end{pmatrix} \in \SL_2(\ZZ)$.
\end{corollary}

\begin{proof}
A lattice $\ZZ\omega_1 + \ZZ\omega_2$ is carried by $z \mapsto z/\omega_1$ to $\ZZ + \ZZ\tau$ with $\tau = \omega_2/\omega_1 \notin \RR$, and replacing $\omega_2$ by $-\omega_2$ if necessary, $\operatorname{Im}\tau > 0$. By
Theorem~\ref{complex-manifolds:theorem-maps-of-tori}, $E_\tau \cong E_{\tau'}$ if and only if $\alpha(\ZZ + \ZZ\tau) = \ZZ + \ZZ\tau'$ for some $\alpha \in \CC^\times$. Then $\alpha\tau = a\tau' + b$ and $\alpha = c\tau' + d$ for an integer matrix,
invertible over $\ZZ$ because the inverse change of basis is also integral; so $\tau = \frac{a\tau' + b}{c\tau' + d}$, and the determinant is $+1$ because both $\tau$ and $\tau'$ lie in the upper half plane, as
$\operatorname{Im}\frac{a\tau' + b}{c\tau' + d} = \frac{(ad - bc)\operatorname{Im}\tau'}{|c\tau' + d|^2}$. Conversely such a matrix gives the lattice isomorphism. The relation is symmetric, so $\tau'$ is also of this form in terms of $\tau$.
\end{proof}

\begin{counterexample}\label{complex-manifolds:counterexample-tori}
Diffeomorphic complex manifolds need not be biholomorphic: all $E_\tau$ are diffeomorphic to $S^1 \times S^1$, but by Corollary~\ref{complex-manifolds:corollary-elliptic-curves} $E_i \not\cong E_{2i}$, since
$\frac{ai + b}{ci + d} = 2i$ with integers of determinant $1$ would give, comparing imaginary parts, $|ci + d|^2 = \frac12$, impossible for integers $c$, $d$.
\end{counterexample}

\begin{proposition}\label{complex-manifolds:proposition-normal-form}
Let $f \colon X \to Y$ be a nonconstant holomorphic map of connected Riemann surfaces (complex manifolds of dimension $1$), and $p \in X$. There are holomorphic coordinates $z$ centred at $p$ and $w$ centred at $f(p)$ in
which $f$ is $z \mapsto z^k$ for a unique integer $k = e_p \geq 1$, the \emph{ramification index}. Consequently $f$ is open, its fibres are discrete, and if $X$ is compact and $Y$ connected, then $f$ is surjective and
has finite fibres.
\end{proposition}

\begin{proof}
Choose coordinates $z$ at $p$ and $w$ at $f(p)$ with $z(p) = 0$, $w(f(p)) = 0$. In them $f$ is a holomorphic function of one variable with $f(0) = 0$, not identically zero near $0$ by the identity theorem
(Theorem~\ref{complex-analysis:theorem-identity}) and connectedness of $X$. So $f(z) = z^kg(z)$ with $k \geq 1$ and $g(0) \neq 0$. On a small disc where $g$ has no zeros, $g'/g$ has a primitive $L$
(Corollary~\ref{complex-analysis:corollary-primitive-convex}), and $(ge^{-L})' = 0$, so $g = ce^L$ with $c \neq 0$; writing $c = e^{c'}$, the function $h = e^{(L + c')/k}$ satisfies $h^k = g$. The new coordinate
$\tilde z = zh(z)$ has derivative $h(0) \neq 0$ at $0$, so it is a local biholomorphism (Theorem~\ref{complex-analysis:theorem-open-mapping}, or
Proposition~\ref{several-complex-variables:proposition-inverse-function}), and $f = \tilde z^k$. The integer $k$ is determined by $f$, as the number of preimages near $p$ of points $w \neq 0$ near $0$. The map
$z \mapsto z^k$ is open with discrete fibres, so $f$ is. If $X$ is compact, $f(X)$ is compact, hence closed, and open, so $f(X) = Y$ when $Y$ is connected; the fibres are discrete and closed in a compact space, hence
finite.
\end{proof}

\begin{proposition}\label{complex-manifolds:proposition-degree}
Let $f \colon X \to Y$ be a nonconstant holomorphic map of compact connected Riemann surfaces. The number $\deg f = \sum_{p \in f^{-1}(q)}e_p$ does not depend on $q \in Y$. It equals the number of preimages of every
point $q$ outside the finite set $f(\{p : e_p > 1\})$ of \emph{branch points}.
\end{proposition}

\begin{proof}
The ramification points, where $f'$ vanishes in coordinates, are the zeros of a holomorphic function locally, hence discrete, hence finitely many in the compact $X$. Fix $q_0$ with preimages $p_1, \ldots, p_r$
(Proposition~\ref{complex-manifolds:proposition-normal-form}), and choose disjoint coordinate discs $D_i$ around $p_i$ on which $f$ is $z \mapsto z^{e_i}$ onto a disc $D$ around $q_0$. The set
$f(X \setminus \bigcup_iD_i)$ is compact and does not contain $q_0$, so there is a neighbourhood $D' \subseteq D$ of $q_0$ avoiding it. For $q \in D' \setminus \{q_0\}$, all preimages of $q$ lie in $\bigcup_iD_i$,
and in $D_i$ the equation $z^{e_i} = w(q) \neq 0$ has exactly $e_i$ solutions, each with ramification index $1$. So the weighted count equals $\sum_ie_i$ on $D'$. Hence $q \mapsto \deg_q f$ is locally constant, hence
constant on the connected $Y$.
\end{proof}

\begin{corollary}\label{complex-manifolds:corollary-maps-to-tori}
Every holomorphic map from a simply connected compact connected complex manifold, such as $\CC\PP^1$, to a complex torus is constant. In particular $\CC\PP^1$ is not biholomorphic, and admits no nonconstant holomorphic
map, to an elliptic curve.
\end{corollary}

\begin{proof}
Let $f \colon X \to T = \CC^m/\Lambda$. The projection $\CC^m \to T$ is a covering (Theorem~\ref{topological-manifolds:theorem-covering-action}), so by the lifting criterion
(Proposition~\ref{homotopy:proposition-lifting-criterion}) $f$ lifts to a continuous $F \colon X \to \CC^m$, which is holomorphic because the projection is a local biholomorphism. By
Proposition~\ref{complex-manifolds:proposition-no-functions} each component of $F$ is constant, so $f$ is constant. The projective line is covered by the two charts $U_0$, $U_1$, each homeomorphic to $\CC$, with
path-connected intersection $\CC^\times$; by the Seifert--van Kampen theorem (Theorem~\ref{homotopy:theorem-van-kampen}) $\pi_1(\CC\PP^1)$ is a quotient of the free product of two trivial groups, hence trivial.
\end{proof}

\section{Meromorphic forms and the residue theorem}\label{complex-manifolds:section-residues}

On a compact Riemann surface, Stokes' theorem turns the residue theorem of complex analysis into a global constraint: the residues of a meromorphic $1$-form add up to zero. This is
the simplest case of the \emph{global residue theorem} \cite[Chapter 5]{griffiths-harris}, and it has quick consequences: a meromorphic function has as many zeros as poles, a complex torus
carries no meromorphic function with a single simple pole, and the roots of a polynomial satisfy identities that are not obvious from the polynomial.

\begin{definition}\label{complex-manifolds:definition-meromorphic-form}
Let $X$ be a Riemann surface. A \emph{meromorphic $1$-form} $\omega$ on $X$ is given in every holomorphic chart $z$ by an expression $f(z)\,dz$ with $f$ meromorphic
(Definition~\ref{complex-analysis:definition-singularity}), such that for two charts $z$ and $w$ the expressions $f(z)\,dz$ and $g(w)\,dw$ agree: $g(w) = f(z(w))\,z^{\prime}(w)$. It is
\emph{holomorphic} if all the $f$ are. The \emph{order} $\operatorname{ord}_p\omega$ is $\operatorname{ord}_0f$ in a chart centred at $p$, and the \emph{residue} $\operatorname{Res}_p\omega$ is the
residue at $0$ of $f$, that is, the coefficient of $z^{-1}$ in its Laurent expansion.
\end{definition}

Since $z^{\prime}(w)$ does not vanish, the order does not depend on the chart; the residue does not either, by the lemma below. For a meromorphic function $h$, $dh = h^{\prime}(z)\,dz$ is a
meromorphic $1$-form, and so is $dh/h$ if $h \not\equiv 0$. On $\CC\PP^1$, with the charts $z$ and $w = 1/z$, the form $dz$ becomes $dz = -dw/w^2$: it has a pole of order $2$ at $\infty$ and residue $0$
there, while $dz/z = -dw/w$ has residue $1$ at $0$ and $-1$ at $\infty$.

Away from its poles, a meromorphic $1$-form is a smooth complex-valued $1$-form, $\omega = \alpha + i\beta$ with real $1$-forms $\alpha$ and $\beta$, and line integrals and integrals over manifolds
of such forms are taken componentwise. It is \emph{closed}: in a chart, with $f$ holomorphic, the Cauchy--Riemann equations (Proposition~\ref{complex-analysis:proposition-cauchy-riemann}) give
$\partial_xf = f^{\prime}$ and $\partial_yf = if^{\prime}$, so $df = f^{\prime}\,(dx + i\,dy) = f^{\prime}\,dz$ and $d(f\,dz) = df \wedge dz = f^{\prime}\,dz \wedge dz = 0$. Holomorphic charts preserve the orientation (the
Jacobian determinant of a holomorphic map is $|f^{\prime}|^2 \geq 0$), so a Riemann surface is canonically oriented, and the boundary of a coordinate disc is traversed counterclockwise in the chart.

\begin{lemma}\label{complex-manifolds:lemma-residue-integral}
Let $\omega$ be a meromorphic $1$-form on a Riemann surface $X$ and $p \in X$. For every closed coordinate disc $D$ around $p$ (the preimage of a closed disc under a holomorphic chart) whose only
possible pole is $p$,
\[
\operatorname{Res}_p\omega = \frac{1}{2\pi i}\int_{\partial D}\omega .
\]
In particular $\operatorname{Res}_p\omega$ does not depend on the chart.
\end{lemma}

\begin{proof}
If $D = \{|z| \leq r\}$ in a chart $z$ centred at $p$ with $\omega = f\,dz$, the line integral of $\omega$ over $\partial D$ is the complex path integral $\int_{|z| = r}f(z)\,dz$, which is
$2\pi i\operatorname{Res}_0f$ by the residue theorem (Theorem~\ref{complex-analysis:theorem-residues}). Now let $D^{\prime}$ be a closed coordinate disc around $p$ for another chart, with no pole other than
$p$. For $r$ small the disc $D_r = \{|z| \leq r\}$ lies in the interior of $D^{\prime}$, and $D^{\prime}$ minus the interior of $D_r$ is a compact manifold with boundary $\partial D^{\prime} \cup \partial D_r$, the
second circle with the opposite orientation. As $\omega$ is closed there, Stokes' theorem (Theorem~\ref{differential-forms:theorem-stokes}) gives $\int_{\partial D^{\prime}}\omega = \int_{\partial D_r}\omega$.
\end{proof}

\begin{theorem}[Residue theorem on a compact Riemann surface]\label{complex-manifolds:theorem-global-residue}
Let $X$ be a compact Riemann surface and $\omega$ a meromorphic $1$-form on $X$. Then $\omega$ has finitely many poles, and
\[
\sum_{p \in X}\operatorname{Res}_p\omega = 0 .
\]
\end{theorem}

\begin{proof}
In a chart the poles of a meromorphic function are isolated, so the poles of $\omega$ form a discrete closed subset of the compact $X$, hence a finite set $\{p_1, \ldots, p_k\}$. Choose pairwise
disjoint closed coordinate discs $D_j$ around the $p_j$. Then $X^{\prime} = X \setminus \bigcup_j\mathrm{int}\,D_j$ is a compact manifold with boundary $\bigcup_j\partial D_j$, each circle with the
orientation opposite to its orientation as the boundary of $D_j$, and $\omega$ is a closed form on $X^{\prime}$. By Stokes' theorem and Lemma~\ref{complex-manifolds:lemma-residue-integral},
\[
0 = \int_{X^{\prime}}d\omega = \int_{\partial X^{\prime}}\omega = -\sum_{j=1}^k\int_{\partial D_j}\omega = -2\pi i\sum_{j=1}^k\operatorname{Res}_{p_j}\omega,
\]
which is the claim.
\end{proof}

The proof is the same as that of the residue theorem in the plane, with one difference: on a compact surface there is no ``outside'', so the integral over the boundary of the complement of the
discs is not the integral over a large circle but zero. For $X = \CC\PP^1$ the point $\infty$ plays the role of the outside: the theorem says that the sum of the residues of a rational form in $\CC$
equals minus its residue at $\infty$.

\begin{corollary}\label{complex-manifolds:corollary-zeros-poles}
Let $f$ be a nonconstant meromorphic function on a compact connected Riemann surface $X$. Then $\sum_p\operatorname{ord}_pf = 0$: $f$ has as many zeros as poles, counted with multiplicity. Hence $f$
takes every value $c \in \CC$ the same number of times, counted with multiplicity, namely the number of its poles.
\end{corollary}

\begin{proof}
At a point $p$, in a chart centred at $p$, write $f = z^kg$ with $k = \operatorname{ord}_pf$ and $g$ holomorphic and nonvanishing near $0$. Then $df/f = k\,dz/z + dg/g$, and $dg/g$ is holomorphic near $0$,
so $\operatorname{Res}_p(df/f) = \operatorname{ord}_pf$. Since $f$ is not identically zero (it is nonconstant on the connected $X$, so not zero on any open set by the identity theorem), $df/f$ is a meromorphic
$1$-form, and Theorem~\ref{complex-manifolds:theorem-global-residue} gives $\sum_p\operatorname{ord}_pf = 0$. Apply this to $f - c$, which has the same poles as $f$.
\end{proof}

This is also Proposition~\ref{complex-manifolds:proposition-degree} applied to $f$ regarded as a holomorphic map to $\CC\PP^1$; the residue theorem gives it without the local analysis of that proof.

\begin{corollary}\label{complex-manifolds:corollary-no-single-pole}
Let $X$ be a compact Riemann surface with a holomorphic $1$-form $\eta$ that does not vanish at $p \in X$. Then no meromorphic function on $X$ is holomorphic on $X \setminus \{p\}$ with a simple pole at $p$.
In particular a meromorphic function on a complex torus $\CC/\Lambda$ that is not constant has at least two poles, counted with multiplicity.
\end{corollary}

\begin{proof}
Suppose $f$ were such a function. In a chart centred at $p$, $f = c/z + (\text{holomorphic})$ with $c \neq 0$ and $\eta = h(z)\,dz$ with $h(0) \neq 0$. Then $f\eta$ is holomorphic except at $p$, where
$\operatorname{Res}_p(f\eta) = c\,h(0) \neq 0$, contradicting Theorem~\ref{complex-manifolds:theorem-global-residue}. On $\CC/\Lambda$ the form $dz$ is invariant under the translations by $\Lambda$, so it
descends to a holomorphic $1$-form without zeros. A nonconstant meromorphic function has a pole, since otherwise it would be a holomorphic function on a compact connected manifold, hence constant
(Proposition~\ref{complex-manifolds:proposition-no-functions}); and a single simple pole is excluded by the first part, applied at that pole.
\end{proof}

On $\CC\PP^1$ the function $z$ has a single simple pole, at $\infty$, so by the corollary $\CC\PP^1$ has no holomorphic $1$-form that is nonzero at $\infty$. In fact it has no nonzero holomorphic
$1$-form at all: such a form is $f(z)\,dz$ with $f$ entire, and $f(z)\,dz = -f(1/w)w^{-2}\,dw$ is holomorphic at $w = 0$ only if $f(z) = O(|z|^{-2})$ as $z \to \infty$; then $f$ is bounded, hence
constant by Liouville's theorem, hence $0$.

\begin{corollary}[Euler--Jacobi formula in one variable]\label{complex-manifolds:corollary-euler-jacobi}
Let $f \in \CC[z]$ have degree $d \geq 2$ and distinct roots $r_1, \ldots, r_d$, and let $g \in \CC[z]$ have degree at most $d - 2$. Then
\[
\sum_{i=1}^d\frac{g(r_i)}{f^{\prime}(r_i)} = 0 .
\]
\end{corollary}

\begin{proof}
Apply Theorem~\ref{complex-manifolds:theorem-global-residue} to $\omega = g(z)\,dz/f(z)$ on $\CC\PP^1$. In $\CC$ its poles are the roots of $f$, which are simple since $f^{\prime}(r_i) \neq 0$, and there
$\operatorname{Res}_{r_i}\omega = \lim_{z \to r_i}(z - r_i)g(z)/f(z) = g(r_i)/f^{\prime}(r_i)$. At $\infty$, with $w = 1/z$, write $g(1/w) = w^{-e}G(w)$ and $f(1/w) = w^{-d}F(w)$ with polynomials $F$, $G$,
where $e = \deg g$ and $F(0) \neq 0$ is the leading coefficient of $f$. Then $\omega = -w^{d-e-2}\frac{G(w)}{F(w)}\,dw$, which is holomorphic at $w = 0$ because $d - e - 2 \geq 0$ (if $g = 0$ there is nothing to prove).
So the residues at the roots add up to zero.
\end{proof}

For example, for $f = z^3 - 1$ the roots are the cube roots of unity $\zeta$, with $f^{\prime}(\zeta) = 3\zeta^2$, and the formula with $g = 1$ and $g = z$ says $\sum_\zeta\zeta^{-2} = 0$ and $\sum_\zeta\zeta^{-1} = 0$;
for $g = z^2$, which has degree $d - 1$, the residue at $\infty$ is $-1$ and indeed $\sum_\zeta\frac{\zeta^2}{3\zeta^2} = 1$. For $f = z^2 - 1$ and $g = 1$: $\frac{1}{2} + \frac{1}{-2} = 0$.

In higher dimension the residue of a meromorphic $n$-form is attached to a point where $n$ divisors meet: if near $P$ the divisors are $\{f_1 = 0\}, \ldots, \{f_n = 0\}$, meeting only at $P$, and
$\omega = g\,dz_1 \wedge \cdots \wedge dz_n/(f_1 \cdots f_n)$ with $g$ holomorphic, the residue is
$\operatorname{Res}_P\omega = \frac{1}{(2\pi i)^n}\int_\Gamma\omega$ over the real $n$-cycle $\Gamma = \{|f_1| = \cdots = |f_n| = \varepsilon\}$, suitably oriented.

\begin{theorem}[Global residue theorem]\label{complex-manifolds:theorem-global-residue-higher}
Let $M$ be a compact complex manifold of dimension $n$, $D_1, \ldots, D_n$ effective divisors whose intersection is a finite set of points, and $\omega$ a meromorphic $n$-form on $M$ whose poles are
contained in $D_1 \cup \cdots \cup D_n$, of order at most one along each $D_i$ (that is, $\omega$ is a holomorphic section of $\Omega^n_M(D_1 + \cdots + D_n)$). Then
$\sum_{P \in D_1 \cap \cdots \cap D_n}\operatorname{Res}_P\omega = 0$.
\end{theorem}

This is \cite[Chapter 5]{griffiths-harris}. For $n = 1$ and $D_1$ the polar divisor it is Theorem~\ref{complex-manifolds:theorem-global-residue} for forms with simple poles. As there, the proof rests on
Stokes' theorem, now applied on the complement of tubes around the divisors, together with a local analysis of the residue. Griffiths and Harris use it to prove the Cayley--Bacharach theorem: if two plane
cubic curves meet in nine distinct points, every cubic through eight of them passes through the ninth.

\section{Uniformisation of Riemann surfaces}\label{complex-manifolds:section-uniformisation}

A \emph{Riemann surface} is a connected complex manifold of dimension $1$. The uniformisation theorem reduces the study of all Riemann surfaces to three simply connected
ones, the sphere, the plane and the disc, and to groups of automorphisms acting on them. Combined with the hyperbolic metric of the disc it shows that most Riemann surfaces,
among them all compact ones of genus at least $2$, carry a natural metric of constant curvature $-1$; and it gives a short proof of Picard's theorem. Throughout,
$\mathbb{D} = \{|z| < 1\}$ and $\mathrm{Aut}(X)$ is the group of biholomorphisms of $X$ to itself.

\begin{lemma}\label{complex-manifolds:lemma-coverings}
Let $X$ be a Riemann surface.
\begin{enumerate}
\item If $p \colon Y \to X$ is a covering map with $Y$ connected, then $Y$ has a unique structure of Riemann surface for which $p$ is a local biholomorphism, and the deck
transformations are biholomorphisms.
\item If a group $\Gamma$ acts on $X$ freely and properly by biholomorphisms, then $X/\Gamma$ has a unique structure of Riemann surface for which $X \to X/\Gamma$ is a local
biholomorphism, and a map $f \colon X/\Gamma \to Z$ to a Riemann surface is holomorphic if and only if its composite with the projection is.
\end{enumerate}
\end{lemma}

\begin{proof}
These are the proofs of Theorem~\ref{smooth-manifolds:theorem-smooth-quotient} and Example~\ref{smooth-manifolds:example-smooth-quotients}, with ``holomorphic'' in place of
``smooth'': the charts of $Y$ (of $X/\Gamma$) are the charts of $X$ composed with $p$ on sheets over evenly covered chart domains (with local inverses of the projection), and their
transition maps are transition maps of $X$, composed with the biholomorphisms of $\Gamma$ in the case of $X/\Gamma$. Topologically, $Y$ is a surface by
Proposition~\ref{topological-manifolds:proposition-constructions}(4), and $X/\Gamma$ one by (3). A deck transformation $\tau$ satisfies $p \circ \tau = p$, so in these charts it is
the identity map of coordinates.
\end{proof}

In particular the universal covering $\widetilde{X}$ of a Riemann surface (Theorem~\ref{homotopy:theorem-universal-cover}) is a simply connected Riemann surface, $\pi_1(X)$ acts
on it freely and properly by biholomorphisms as the deck group (Proposition~\ref{homotopy:proposition-deck}), and $X \cong \widetilde{X}/\pi_1(X)$.

\begin{theorem}[Uniformisation theorem]\label{complex-manifolds:theorem-uniformisation}
Every simply connected Riemann surface is biholomorphic to exactly one of $\CC\PP^1$, $\CC$ and $\mathbb{D}$.
\end{theorem}

This is due to Koebe and Poincar\'e (1907); see \cite{forster-riemann-surfaces}. The three surfaces are pairwise non-biholomorphic: $\CC\PP^1$ is compact and the others are not,
and a biholomorphism $\CC \to \mathbb{D}$ would be a bounded nonconstant entire function (Exercise~\ref{complex-analysis:exercise-disc-not-plane}). For open subsets of $\CC$ the
theorem is the Riemann mapping theorem (Theorem~\ref{complex-analysis:theorem-riemann-mapping}); the general case constructs the biholomorphism from Green's functions, using
Perron's method to solve the Dirichlet problem on the surface.

\begin{lemma}\label{complex-manifolds:lemma-automorphisms}
\begin{enumerate}
\item $\mathrm{Aut}(\CC\PP^1)$ consists of the M\"obius maps $z \mapsto (az + b)/(cz + d)$, $ad - bc \neq 0$, and each of them has a fixed point.
\item $\mathrm{Aut}(\CC)$ consists of the maps $z \mapsto az + b$, $a \neq 0$; those without fixed points are the translations $z \mapsto z + b$, $b \neq 0$.
\item $\mathrm{Aut}(\mathbb{D})$ consists of the maps $z \mapsto \lambda(z - a)/(1 - \bar az)$ with $|\lambda| = 1$ and $a \in \mathbb{D}$, and acts transitively on $\mathbb{D}$.
\end{enumerate}
\end{lemma}

\begin{proof}
(2) is Exercise~\ref{complex-analysis:exercise-automorphisms-plane}, and $az + b = z$ has the solution $b/(1 - a)$ when $a \neq 1$. (3) is
Proposition~\ref{complex-analysis:proposition-disc-automorphisms}; the map $\phi_{-a}$ sends $0$ to $a$. (1) A M\"obius map is holomorphic in the charts of $\CC\PP^1$, with the
M\"obius inverse $z \mapsto (dz - b)/(-cz + a)$, so it is an automorphism, and the M\"obius maps act transitively ($z \mapsto z + b$ and $z \mapsto 1/z$ move any point to any
other). Given $f \in \mathrm{Aut}(\CC\PP^1)$, choose a M\"obius $m$ with $m(f(\infty)) = \infty$; then $m \circ f$ restricts to an automorphism of $\CC = \CC\PP^1 \setminus \{\infty\}$,
hence is $z \mapsto az + b$ by (2), and $f = m^{-1} \circ (m \circ f)$ is M\"obius. If $c = 0$, $\infty$ is a fixed point; if $c \neq 0$, the equation $cz^2 + (d - a)z - b = 0$ has a
root $z_0 \in \CC$ (Corollary~\ref{complex-analysis:corollary-liouville}(2)), and $cz_0 + d \neq 0$, since at $z_0 = -d/c$ the left side is $(ad - bc)/c \neq 0$; so $z_0$ is a fixed point.
\end{proof}

\begin{proposition}[The Poincar\'e metric]\label{complex-manifolds:proposition-poincare-metric}
Let $g_{\mathbb{D}} = 4|dz|^2/(1 - |z|^2)^2$ on $\mathbb{D}$, that is $4(dx^2 + dy^2)/(1 - x^2 - y^2)^2$.
\begin{enumerate}
\item (Schwarz--Pick) For every holomorphic $f \colon \mathbb{D} \to \mathbb{D}$,
\[
\frac{|f^{\prime}(z)|}{1 - |f(z)|^2} \leq \frac1{1 - |z|^2},
\]
that is, $f^*g_{\mathbb{D}} \leq g_{\mathbb{D}}$: holomorphic self-maps of the disc do not increase lengths of curves, hence distances.
\item Every automorphism of $\mathbb{D}$ is an isometry of $g_{\mathbb{D}}$.
\item $g_{\mathbb{D}}$ is complete, with constant curvature $-1$.
\end{enumerate}
\end{proposition}

\begin{proof}
(1) Fix $a \in \mathbb{D}$, let $b = f(a)$ and $\phi_a(z) = (z - a)/(1 - \bar az)$. A computation gives $\phi_a^{\prime}(z) = (1 - |a|^2)/(1 - \bar az)^2$, so
$\phi_{-a}^{\prime}(0) = 1 - |a|^2$ and $\phi_b^{\prime}(b) = 1/(1 - |b|^2)$. The map $h = \phi_b \circ f \circ \phi_{-a}$ sends $\mathbb{D}$ to $\mathbb{D}$ and $0$ to $0$, so
$|h^{\prime}(0)| \leq 1$ by Schwarz's lemma (Lemma~\ref{complex-analysis:lemma-schwarz}), and by the chain rule $h^{\prime}(0) = f^{\prime}(a)(1 - |a|^2)/(1 - |b|^2)$. At a point
$z$, $f^*g_{\mathbb{D}}$ is $4|f^{\prime}(z)|^2|dz|^2/(1 - |f(z)|^2)^2$, so the inequality says $f^*g_{\mathbb{D}} \leq g_{\mathbb{D}}$, and lengths of curves decrease.

(2) Apply (1) to an automorphism and to its inverse: $f^*g_{\mathbb{D}} \leq g_{\mathbb{D}}$ and $g_{\mathbb{D}} = (f^{-1})^*f^*g_{\mathbb{D}}$ force equality.

(3) The curvature is $-1$ by Proposition~\ref{riemannian-geometry:proposition-conformal-curvature} (Example~\ref{riemannian-geometry:example-conformal-metrics}(2)). For
completeness, choose $\varepsilon > 0$ such that $\exp_0$ is defined on the ball of radius $\varepsilon$ in $T_0\mathbb{D}$. Since $\mathrm{Aut}(\mathbb{D})$ acts transitively
by isometries, and isometries map geodesics to geodesics, the same $\varepsilon$ works at every point: every geodesic of unit speed is defined on an interval of length
$\varepsilon$ around each of its times. So a maximal unit speed geodesic defined on $[0, T)$ with $T < \infty$ could be continued beyond $T$ from its point at time
$T - \varepsilon/2$, a contradiction. Geodesics are defined for all time, and $g_{\mathbb{D}}$ is complete by Theorem~\ref{riemannian-geometry:theorem-hopf-rinow}.
\end{proof}

\begin{theorem}\label{complex-manifolds:theorem-riemann-surfaces-classification}
Every Riemann surface $X$ is biholomorphic to $\widetilde{X}/\Gamma$, where $\widetilde{X} \in \{\CC\PP^1, \CC, \mathbb{D}\}$ is its universal covering and $\Gamma \cong \pi_1(X)$ is
a group of automorphisms of $\widetilde{X}$ acting freely and properly. Moreover:
\begin{enumerate}
\item $\widetilde{X} = \CC\PP^1$ if and only if $X \cong \CC\PP^1$;
\item $\widetilde{X} = \CC$ if and only if $X$ is biholomorphic to $\CC$, to $\CC^\times$ or to a complex torus $\CC/\Lambda$;
\item in all other cases $\widetilde{X} = \mathbb{D}$, and $X$ carries a unique complete metric of constant curvature $-1$ for which the projection $\mathbb{D} \to X$ is a local
isometry; in holomorphic coordinates it has the form $\lambda(z)|dz|^2$.
\end{enumerate}
\end{theorem}

\begin{proof}
The first statement is Theorem~\ref{complex-manifolds:theorem-uniformisation} and the remark after Lemma~\ref{complex-manifolds:lemma-coverings}. (1) By
Lemma~\ref{complex-manifolds:lemma-automorphisms}(1) every automorphism of $\CC\PP^1$ has a fixed point, so a group acting freely on it is trivial. (2) By
Lemma~\ref{complex-manifolds:lemma-automorphisms}(2), $\Gamma$ consists of translations $z \mapsto z + \lambda$, and $\Lambda = \{\lambda\}$ is a subgroup of $\CC = \RR^2$. It is
discrete: it is the orbit of $0$, which is a fibre of the covering $\CC \to X$. By Lemma~\ref{lie-groups:lemma-discrete-subgroups}, $\Lambda$ is $0$, $\ZZ\omega$ with
$\omega \neq 0$, or a lattice. In the first case $X \cong \CC$; in the second, $z \mapsto e^{2\pi iz/\omega}$ is a covering $\CC \to \CC^\times$ whose fibres are the cosets of
$\ZZ\omega$, so it induces a bijective local biholomorphism $\CC/\ZZ\omega \to \CC^\times$; in the third $X$ is a complex torus. Conversely these surfaces have universal covering
$\CC$. (3) The group $\Gamma \subseteq \mathrm{Aut}(\mathbb{D})$ acts by isometries of $g_{\mathbb{D}}$ (Proposition~\ref{complex-manifolds:proposition-poincare-metric}), so
$g_{\mathbb{D}}$ descends to a metric on $X$ making the projection a local isometry (Proposition~\ref{riemannian-geometry:proposition-quotient-metric}), of curvature $-1$; it is
complete because geodesics of $X$ lift to geodesics of the complete $\mathbb{D}$ (as in the proof of Theorem~\ref{riemannian-geometry:theorem-bonnet-myers}), and in the
holomorphic coordinates given by local inverses of the projection it is $4|dz|^2/(1 - |z|^2)^2$. It is unique because the projection is a local diffeomorphism, so the metric on
$X$ is determined by its pullback $g_{\mathbb{D}}$.
\end{proof}

\begin{corollary}\label{complex-manifolds:corollary-compact-uniformisation}
Let $X$ be a compact Riemann surface of genus $g$ (Theorem~\ref{topological-manifolds:theorem-classification-surfaces}; $X$ is orientable). If $g = 0$ then $X \cong \CC\PP^1$; if
$g = 1$ then $X$ is a complex torus $\CC/\Lambda$; and if $g \geq 2$ then $X \cong \mathbb{D}/\Gamma$, and $X$ carries a metric of constant curvature $-1$, compatible with its
complex structure, with area $4\pi(g - 1)$.
\end{corollary}

\begin{proof}
If $g = 0$, $X$ is simply connected (Corollary~\ref{homotopy:corollary-spheres}), so $X$ is its own universal covering, compact, hence $\CC\PP^1$. If $g \geq 1$, $\pi_1(X)$ is
nontrivial, as its abelianisation $H_1(X; \ZZ)$ (Theorem~\ref{singular-homology:theorem-hurewicz}) is $\ZZ^{2g}$
(Example~\ref{singular-homology:example-cellular-computations}), so $\widetilde{X} \neq \CC\PP^1$ by
Theorem~\ref{complex-manifolds:theorem-riemann-surfaces-classification}(1). If $\widetilde{X} = \mathbb{D}$, the metric of curvature $-1$ has area $A$ with $-A = 2\pi(2 - 2g)$ by
the Gauss--Bonnet theorem (Theorem~\ref{riemannian-geometry:theorem-gauss-bonnet}), so $A = 4\pi(g - 1)$ and $g \geq 2$. If $\widetilde{X} = \CC$, then $X$ is compact among
$\CC$, $\CC^\times$ and the tori, so it is a torus, of genus $1$. Hence $g = 1$ forces $\widetilde{X} = \CC$ and $g \geq 2$ forces $\widetilde{X} = \mathbb{D}$.
\end{proof}

\begin{corollary}[Little Picard theorem]\label{complex-manifolds:corollary-picard}
A nonconstant entire function takes every complex value with at most one exception.
\end{corollary}

\begin{proof}
Suppose the entire function $f$ omits two values; composing with an affine map we may assume they are $0$ and $1$, so $f \colon \CC \to X = \CC \setminus \{0, 1\}$. The fundamental
group of $X$ is not abelian: $X = U \cup V$ with $U = \{\operatorname{Re}z < \frac34\} \setminus \{0\}$ and $V = \{\operatorname{Re}z > \frac14\} \setminus \{1\}$, where
$U \cap V$ is a convex strip, hence simply connected, and $U$, $V$ deformation retract onto the circles of radius $\frac14$ around $0$ and $1$ (radially, along segments
that stay in the convex half planes), so $\pi_1(U) \cong \pi_1(V) \cong \ZZ$ (Lemma~\ref{homotopy:lemma-deformation-retract} and Theorem~\ref{homotopy:theorem-circle}) and
$\pi_1(X) \cong \ZZ * \ZZ$ by Corollary~\ref{homotopy:corollary-spheres}, a non-abelian group (Example~\ref{homotopy:example-wedge}). Since $X$ is not compact and $\pi_1(X)$ is not abelian, while
$\CC$, $\CC^\times$ and the tori have abelian fundamental groups $0$, $\ZZ$ and $\ZZ^2$, Theorem~\ref{complex-manifolds:theorem-riemann-surfaces-classification} gives a covering
map $p \colon \mathbb{D} \to X$ that is a local biholomorphism. As $\CC$ is simply connected, $f$ lifts to a continuous $\tilde f \colon \CC \to \mathbb{D}$ with $p \circ \tilde f = f$
(Proposition~\ref{homotopy:proposition-lifting-criterion}), which is holomorphic, being locally a local inverse of $p$ composed with $f$. Then $\tilde f$ is a bounded entire function,
hence constant (Corollary~\ref{complex-analysis:corollary-liouville}), and so is $f$.
\end{proof}

\begin{remark}\label{complex-manifolds:remark-uniformisation-metrics}
For compact surfaces the three cases match the sign of the Euler characteristic: $\CC\PP^1$ carries the round metric of curvature $1$ (Example~\ref{riemannian-geometry:example-conformal-metrics}(3)),
complex tori the flat metric $|dz|^2$, and surfaces of genus $g \geq 2$ the hyperbolic metric, all compatible with the complex structure; the sign of the curvature is forced by the
Gauss--Bonnet theorem. Schwarz--Pick also shows that the hyperbolic metric is intrinsic: a holomorphic map $f \colon X \to Y$ between Riemann surfaces covered by the disc
lifts to a holomorphic map $\mathbb{D} \to \mathbb{D}$ of the universal coverings (Proposition~\ref{homotopy:proposition-lifting-criterion}), so it does not increase the lengths of
curves for the hyperbolic metrics of Theorem~\ref{complex-manifolds:theorem-riemann-surfaces-classification}(3).
\end{remark}

\section{Almost complex structures and forms of type $(p,q)$}\label{complex-manifolds:section-pq}

\begin{definition}\label{complex-manifolds:definition-almost-complex}
An \emph{almost complex structure} on a smooth manifold $M$ is an endomorphism $J$ of $TM$ with $J^2 = -\id$. It makes each $T_pM$ a complex vector space, and gives the decomposition
$T_pM \otimes \CC = T^{1,0}_p \oplus T^{0,1}_p$ into the $\pm i$ eigenspaces, with $T^{0,1} = \overline{T^{1,0}}$, and dually
$\Lambda^k T^*_pM \otimes \CC = \bigoplus_{p+q=k}\Lambda^{p,q}$ (Proposition~\ref{multilinear-algebra:proposition-pq-decomposition}). A complex manifold has a canonical almost complex structure, given in
holomorphic coordinates $z^j = x^j + iy^j$ by $J\partial_{x^j} = \partial_{y^j}$, $J\partial_{y^j} = -\partial_{x^j}$; then $T^{1,0}$ has frame $\partial/\partial z^j$ and $\Lambda^{p,q}$ has frame
$dz^I \wedge d\bar z^J$ with $|I| = p$, $|J| = q$. An almost complex structure is \emph{integrable} if it arises from a complex structure.
\end{definition}

\begin{definition}\label{complex-manifolds:definition-dbar}
On a complex manifold, $\Omega^{p,q}(X)$ denotes the smooth $(p,q)$-forms, so $\Omega^k(X, \CC) = \bigoplus_{p+q=k}\Omega^{p,q}(X)$. Writing $d = \partial + \bar\partial$ with
$\partial \colon \Omega^{p,q} \to \Omega^{p+1,q}$ and $\bar\partial \colon \Omega^{p,q} \to \Omega^{p,q+1}$ the components of $d$, one has in coordinates
\[
\partial\left(\sum f_{IJ}\,dz^I \wedge d\bar z^J\right) = \sum\frac{\partial f_{IJ}}{\partial z^k}\,dz^k \wedge dz^I \wedge d\bar z^J, \qquad
\bar\partial\left(\sum f_{IJ}\,dz^I \wedge d\bar z^J\right) = \sum\frac{\partial f_{IJ}}{\partial\bar z^k}\,d\bar z^k \wedge dz^I \wedge d\bar z^J .
\]
\end{definition}

\begin{proposition}\label{complex-manifolds:proposition-dbar}
On a complex manifold $d = \partial + \bar\partial$, and $\partial^2 = \bar\partial^2 = \partial\bar\partial + \bar\partial\partial = 0$. A function is holomorphic if and only if $\bar\partial f = 0$, and a $(p, 0)$-form
$\alpha$ is \emph{holomorphic} (has holomorphic coefficients) if and only if $\bar\partial\alpha = 0$; the sheaf of holomorphic $p$-forms is denoted $\Omega^p_X$, with $\Omega^0_X = \mathcal{O}_X$ and
$\Omega^n_X = K_X$ the \emph{canonical bundle}.
\end{proposition}

\begin{proof}
In holomorphic coordinates $d$ of a $(p,q)$-form has components of types $(p+1, q)$ and $(p, q+1)$ only, by the coordinate formula for $d$; these are $\partial$ and $\bar\partial$. Then $0 = d^2$ decomposes into
its components of types $(p+2, q)$, $(p+1, q+1)$ and $(p, q+2)$, which are $\partial^2$, $\partial\bar\partial + \bar\partial\partial$ and $\bar\partial^2$. The statement about functions is the Cauchy--Riemann
equations (Definition~\ref{several-complex-variables:definition-holomorphic}), and for forms it is the same condition on the coefficients.
\end{proof}

\begin{theorem}[Newlander--Nirenberg]\label{complex-manifolds:theorem-newlander-nirenberg}
An almost complex structure $J$ on a smooth manifold is integrable if and only if the \emph{Nijenhuis tensor} vanishes,
\[
N_J(X, Y) = [JX, JY] - J[JX, Y] - J[X, JY] - [X, Y] = 0,
\]
equivalently if and only if $T^{0,1}$ is closed under the Lie bracket, equivalently if and only if $d = \partial + \bar\partial$ on $1$-forms with respect to the induced type decomposition.
\end{theorem}

The implication ``integrable $\Rightarrow$ $N_J = 0$'' is a computation in holomorphic coordinates, where the brackets of the coordinate fields vanish; the converse is a hard analytic theorem, proved in
\cite[Chapter 2]{huybrechts-complex} and, in the real-analytic case, by the Frobenius theorem (Theorem~\ref{smooth-manifolds:theorem-frobenius}) applied to $T^{0,1}$ after complexification.

\begin{proposition}\label{complex-manifolds:proposition-surfaces-almost-complex}
Every almost complex structure on a manifold of real dimension $2$ has vanishing Nijenhuis tensor. Consequently, by Theorem~\ref{complex-manifolds:theorem-newlander-nirenberg}, every oriented surface with a Riemannian metric is a Riemann
surface, with $J$ the rotation by $\pi/2$ in each tangent plane.
\end{proposition}

\begin{proof}
$N_J$ is antisymmetric and $C^\infty$-bilinear, so it suffices to evaluate it on $X$ and $JX$ for a local nowhere vanishing field $X$, which together with $JX$ form a frame. Using $J^2 = -1$,
$N_J(X, JX) = [JX, -X] - J[JX, JX] - J[X, -X] - [X, JX] = [X, JX] - [X, JX] = 0$. On an oriented surface with a metric, rotation by $\pi/2$ in the positive sense is a smooth endomorphism with $J^2 = -1$.
\end{proof}

\section{Analytic subvarieties and blow-ups}\label{complex-manifolds:section-subvarieties}

\begin{definition}\label{complex-manifolds:definition-analytic-subvariety}
An \emph{analytic subvariety} of a complex manifold $X$ is a closed subset that is locally the common zero set of finitely many holomorphic functions (Definition~\ref{several-complex-variables:definition-analytic-set});
it is \emph{smooth}, that is, a complex submanifold, at its regular points. By Theorem~\ref{several-complex-variables:theorem-analytic-structure} the singular locus is a proper analytic subvariety, the regular
locus is dense, and $Z$ has a well-defined dimension on each irreducible component.
\end{definition}

\begin{proposition}\label{complex-manifolds:proposition-subvariety-class}
Let $Z \subseteq X$ be an analytic subvariety of pure complex dimension $k$ in a compact complex manifold of dimension $n$. Then integration over the regular locus,
$\varphi \mapsto \int_{Z_{\mathrm{reg}}}\varphi$, converges for every smooth $2k$-form $\varphi$ and defines a closed current of degree $2n - 2k$; its class in $H^{2n-2k}(X; \RR)$ is the \emph{cycle class} $[Z]$.
Moreover $\int_Z\varphi \geq 0$ when $\varphi$ is a positive $(k,k)$-form.
\end{proposition}

\begin{proof}[Sketch of proof]
Convergence follows from the local structure of analytic sets: near a regular point $Z$ is a submanifold, and near a singular point it is a finite branched cover of a polydisc
(Theorem~\ref{several-complex-variables:theorem-analytic-structure}), so the volume of $Z \cap K$ is finite for compact $K$ (in a chart, Wirtinger's inequality bounds the volume by the integral of a power of the
standard form $\frac{i}{2}\sum_jdz^j \wedge d\bar z^j$ of $\CC^n$). Closedness, that is, Stokes' theorem across the singular locus, is proved by cutting out a neighbourhood of the singular set, whose boundary
has small area because the singular locus has real codimension at least $2$. The positivity statement is local linear algebra. Details are in \cite[Chapter 0, Section 2]{griffiths-harris} and
\cite[Chapter 11]{voisin-hodge-1}.
\end{proof}

\begin{remark}\label{complex-manifolds:remark-cycle-class-nonzero}
If $X$ carries a K\"ahler form $\omega$ (Chapter~\ref{kahler}) and $Z \neq \emptyset$, then $[Z] \neq 0$: pairing with the closed form $\omega^k$ gives $\int_Z\omega^k > 0$, the volume of $Z$ up to a constant.
On non-K\"ahler manifolds cycle classes can vanish; for instance on the Hopf surface of Exercise~\ref{complex-manifolds:exercise-hopf-surface} the elliptic curves $\{z_1 = 0\}/\langle 2 \rangle$ have
zero class, since $H^2$ of the Hopf surface vanishes.
\end{remark}

\begin{definition}\label{complex-manifolds:definition-blowup}
The \emph{blow-up} of $\CC^n$ at the origin is
\[
\mathrm{Bl}_0\CC^n = \{(z, \ell) \in \CC^n \times \CC\PP^{n-1} : z \in \ell\},
\]
with the projection $\pi$ to $\CC^n$. It is a complex manifold, $\pi$ is a biholomorphism over $\CC^n \setminus \{0\}$, and $E = \pi^{-1}(0) \cong \CC\PP^{n-1}$ is the \emph{exceptional divisor}. For a point $p$ of
a complex manifold $X$ one defines $\mathrm{Bl}_pX$ by gluing $\mathrm{Bl}_0\CC^n$ in a chart, and similarly along a complex submanifold $Z \subseteq X$, with $E = \PP(N_{Z/X})$.
\end{definition}

\begin{proposition}\label{complex-manifolds:proposition-blowup}
$\mathrm{Bl}_pX$ is a compact complex manifold if $X$ is, $\pi$ is proper and holomorphic, and
\[
H^k(\mathrm{Bl}_pX; \ZZ) \cong H^k(X; \ZZ) \oplus \bigoplus_{0 < j < n} H^{k-2j}(\mathrm{pt}; \ZZ),
\]
so blowing up a point adds one generator in each even degree $2, 4, \ldots, 2n-2$. The exceptional divisor has normal bundle $\mathcal{O}(-1)$ in $\mathrm{Bl}_0\CC^n$.
\end{proposition}

\begin{proof}
In the chart $\ell = [1 : w_2 : \cdots : w_n]$ the map $(u, w_2, \ldots, w_n) \mapsto (u, uw_2, \ldots, uw_n)$ identifies $\CC^n$ with an open subset of $\mathrm{Bl}_0\CC^n$, which shows that the blow-up is a complex
manifold and that $\pi$ is holomorphic and biholomorphic away from $E$; properness holds because $\CC\PP^{n-1}$ is compact. In these coordinates $E$ is $\{u = 0\}$, and the transition functions of its
normal bundle are those of the tautological bundle, so $N_{E} \cong \mathcal{O}(-1)$. For the cohomology, use the Mayer--Vietoris sequence for $\mathrm{Bl}_pX = (\mathrm{Bl}_pX \setminus E) \cup \pi^{-1}(B)$ with
$B$ a small ball, together with $\pi^{-1}(B) \simeq \CC\PP^{n-1}$ and $\mathrm{Bl}_pX \setminus E \cong X \setminus \{p\}$, and Example~\ref{singular-homology:example-cohomology-groups}.
\end{proof}

\begin{remark}\label{complex-manifolds:remark-blowup-birational}
Blow-ups are the basic birational modifications: Hironaka's resolution of singularities states that every analytic (or algebraic, in characteristic zero) variety becomes smooth after finitely many blow-ups
along smooth centres, and the weak factorisation theorem says that any two birational smooth projective varieties are connected by blow-ups and blow-downs. Both are used to reduce statements about the Hodge
conjecture to the smooth projective case (Chapter~\ref{mixed-hodge-structures}).
\end{remark}

\section{Holomorphic Vector Bundles}\label{complex-manifolds:section-holomorphic-bundles}
Holomorphic line bundles are the objects that encode divisors, projective embeddings and, through their first Chern classes, the
$(1,1)$-part of cohomology. The exponential sequence (Theorem~\ref{complex-manifolds:theorem-exponential}) computes the Picard
group in terms of cohomology and already contains the case $p = 1$ of the Hodge conjecture: the classes of divisors are exactly the
integral classes of type $(1,1)$ (Chapter~\ref{hodge-conjecture}). References: \cite[Chapter 1]{griffiths-harris},
\cite[Chapters 2--4]{huybrechts-complex}, \cite[Chapter 4]{voisin-hodge-1}.

\noindent\textbf{Prerequisites.} Chapter~\ref{complex-manifolds} (Complex Manifolds), Section~\ref{sheaves:section-sheaf-cohomology} (Sheaf
Cohomology) and Chapter~\ref{characteristic-classes} (Characteristic Classes).

$X$ denotes a complex manifold of dimension $n$.

\section{Holomorphic bundles and sheaves}\label{complex-manifolds:section-bundles}

\begin{definition}\label{complex-manifolds:definition-holomorphic-bundle}
A \emph{holomorphic vector bundle} of rank $r$ on $X$ is a complex manifold $E$ with a holomorphic map $\pi \colon E \to X$ and $\CC$-vector space structures on the fibres, locally biholomorphic to
$U \times \CC^r$ over $U$ and linear on fibres; equivalently, it is given by holomorphic transition functions $g_{ij} \colon U_i \cap U_j \to \GL_r(\CC)$ satisfying the cocycle condition
(Proposition~\ref{vector-bundles:proposition-cocycles}). Its sheaf of holomorphic sections is a locally free $\mathcal{O}_X$-module of rank $r$, and this is an equivalence between holomorphic bundles and
locally free sheaves, as in Proposition~\ref{vector-bundles:proposition-sheaf-of-sections}.
\end{definition}

\begin{example}\label{complex-manifolds:example-bundles}
The holomorphic tangent bundle $T^{1,0}X$, with transition functions the holomorphic Jacobians, and its dual, the holomorphic cotangent bundle
(Section~\ref{complex-manifolds:section-tangent}); the bundles $\Omega^p_X = \Lambda^p(T^{1,0}X)^\vee$ of holomorphic $p$-forms and the canonical bundle
$K_X = \Omega^n_X$; the tautological bundle $\mathcal{O}(-1)$ on $\CC\PP^n$ (Exercise~\ref{vector-bundles:exercise-tautological}), its dual $\mathcal{O}(1)$ and the powers
$\mathcal{O}(k) = \mathcal{O}(1)^{\otimes k}$. The holomorphic sections of $\mathcal{O}(k)$ on $\CC\PP^n$ are the homogeneous polynomials of degree $k$ for $k \geq 0$, and $0$ for $k < 0$
(Exercise~\ref{complex-manifolds:exercise-sections-ok}).
\end{example}

\begin{proposition}\label{complex-manifolds:proposition-picard}
The isomorphism classes of holomorphic line bundles on $X$ form a group $\Pic(X)$ under $\otimes$, and
\[
\Pic(X) \cong H^1(X, \mathcal{O}_X^*),
\]
where $\mathcal{O}_X^*$ is the sheaf of nowhere vanishing holomorphic functions.
\end{proposition}

\begin{proof}
As in Proposition~\ref{vector-bundles:proposition-cocycles} and Theorem~\ref{vector-bundles:theorem-line-bundles}, line bundles correspond to \v{C}ech cocycles with values in $\mathcal{O}_X^*$ modulo
coboundaries, and $\check{H}^1 = H^1$ (Remark~\ref{sheaves:remark-cech-limit}); tensor product corresponds to multiplication of cocycles.
\end{proof}

\begin{theorem}[Exponential sequence]\label{complex-manifolds:theorem-exponential}
The sequence of sheaves $0 \to \ZZ_X \to \mathcal{O}_X \xrightarrow{\exp(2\pi i\,\cdot)} \mathcal{O}_X^* \to 0$ is exact, and its long exact sequence reads
\[
\cdots \to H^1(X, \ZZ) \to H^1(X, \mathcal{O}_X) \to \Pic(X) \xrightarrow{\ c_1\ } H^2(X, \ZZ) \to H^2(X, \mathcal{O}_X) \to \cdots
\]
The map $c_1$ is the first Chern class of the underlying smooth line bundle (Theorem~\ref{vector-bundles:theorem-line-bundles}). In particular, for compact $X$:
\begin{enumerate}
\item $c_1$ is surjective onto the kernel of $H^2(X, \ZZ) \to H^2(X, \mathcal{O}_X)$;
\item its kernel $\Pic^0(X)$ is $H^1(X, \mathcal{O}_X)/\operatorname{im}H^1(X, \ZZ)$;
\item if $H^1(X, \mathcal{O}_X) = H^2(X, \mathcal{O}_X) = 0$, then $c_1 \colon \Pic(X) \to H^2(X, \ZZ)$ is an isomorphism.
\end{enumerate}
\end{theorem}

\begin{proof}
Exactness: the kernel of $\exp(2\pi i \cdot)$ is the constant sheaf $\ZZ$, and surjectivity holds on stalks because a nonvanishing holomorphic function on a small ball has a holomorphic logarithm (integrate
$df/f$, which is possible on a simply connected domain by Theorem~\ref{complex-analysis:theorem-cauchy-homotopy} applied in one variable at a time, or note that the ball is contractible and use the same
argument as in Theorem~\ref{vector-bundles:theorem-line-bundles}). The long exact sequence is Theorem~\ref{homological-algebra:theorem-long-exact-sequence}, using
$\Pic(X) \cong H^1(X, \mathcal{O}_X^*)$. That the connecting map is the topological first Chern class follows by comparing with the corresponding sequence for smooth functions, into which this one maps, and
using that the sheaf of smooth functions is acyclic. The three consequences are exactness of the long sequence.
\end{proof}

\begin{remark}\label{complex-manifolds:remark-picard-variety}
When $X$ is compact K\"ahler, the Hodge decomposition (Chapter~\ref{kahler}) shows that $H^1(X, \ZZ)$ maps onto a lattice in the complex vector space $H^1(X, \mathcal{O}_X)$ of dimension
$b_1/2$, so $\Pic^0(X)$ is a compact complex torus, the \emph{Picard variety}, and the image of $c_1$, the \emph{N\'eron--Severi group}, is the group of integral classes of type $(1,1)$
(Chapter~\ref{hodge-conjecture}).
\end{remark}

\begin{lemma}\label{complex-manifolds:lemma-trivial-sections}
Let $L$ be a holomorphic line bundle on $X$.
\begin{enumerate}
\item $L$ is trivial if and only if it has a nowhere vanishing holomorphic section.
\item If $X$ is compact and connected and both $L$ and its dual $L^\vee$ have nonzero holomorphic sections, then $L$ is trivial. In particular, if $L$ is nontrivial and $H^0(X, L) \neq 0$, then
$H^0(X, L^\vee) = 0$.
\end{enumerate}
\end{lemma}

\begin{proof}
(1) A nowhere vanishing section $s$ gives the map $\mathcal{O}_X \to L$, $f \mapsto fs$, which is an isomorphism on every fibre; conversely the trivial bundle has the section $1$.

(2) First, if $U$ is a connected open set and $a$, $b$ are holomorphic functions on $U$, neither identically zero, then $ab$ is not identically zero: otherwise $U \setminus a^{-1}(0)$ is open, nonempty by the identity
theorem (Proposition~\ref{several-complex-variables:proposition-identity-maximum}), and $b$ vanishes on it, so $b = 0$ by the identity theorem again. Next, a nonzero section $s$ of $L$ does not vanish identically on any
nonempty open set: the set of points near which $s$ vanishes identically is open, and it is closed by the identity theorem applied in local frames, so it is empty because $X$ is connected and $s \neq 0$.

Now let $s \in H^0(X, L)$ and $t \in H^0(X, L^\vee)$ be nonzero. The pairing $t(s)$ is a holomorphic function on $X$, hence a constant $c$ (Proposition~\ref{complex-manifolds:proposition-no-functions}). In a
connected coordinate neighbourhood $U$ with a holomorphic frame $e$ and dual frame $e^\vee$, write $s = ae$ and $t = be^\vee$; then $t(s) = ab$ with $a$, $b$ not identically zero on $U$, so $c \neq 0$. Therefore
$a(x)b(x) = c \neq 0$ everywhere, $s$ vanishes nowhere, and $L$ is trivial by (1).
\end{proof}

\begin{definition}\label{complex-manifolds:definition-dbar-operator}
Let $L$ be a smooth complex line bundle on $X$, and let $\mathcal{A}^{0,q}(L)$ denote the smooth $L$-valued $(0,q)$-forms. A \emph{$\bar\partial$-operator} on $L$ is a $\CC$-linear map
$D \colon \mathcal{A}^0(L) \to \mathcal{A}^{0,1}(L)$ with $D(fs) = \bar\partial f \otimes s + fDs$ for smooth functions $f$. In a local smooth frame $e$ it has the form $D(fe) = (\bar\partial f + fa_e)e$ for a
$(0,1)$-form $a_e = De/e$. The operator is \emph{integrable} if $\bar\partial a_e = 0$ for every local frame $e$.
\end{definition}

\begin{proposition}\label{complex-manifolds:proposition-dbar-operator}
\begin{enumerate}
\item If $L$ is a holomorphic line bundle, the operator $\bar\partial_L(fe) = \bar\partial f \otimes e$, for $e$ a holomorphic frame, is a well-defined integrable $\bar\partial$-operator, and a smooth section $s$ is
holomorphic if and only if $\bar\partial_Ls = 0$.
\item Integrability does not depend on the frame: if $e' = ge$, then $a_{e'} = a_e + g^{-1}\bar\partial g$ and $\bar\partial a_{e'} = \bar\partial a_e$.
\item Conversely, for every integrable $\bar\partial$-operator $D$ on a smooth complex line bundle $L$ there is a unique structure of holomorphic line bundle on $L$ with $D = \bar\partial_L$.
\end{enumerate}
In particular, on a Riemann surface, where there are no $(0,2)$-forms, every $\bar\partial$-operator is integrable.
\end{proposition}

\begin{proof}
(2) $D(e') = D(ge) = (\bar\partial g + ga_e)e = (g^{-1}\bar\partial g + a_e)e'$. Moreover $\bar\partial(g^{-1}\bar\partial g) = -g^{-2}\bar\partial g \wedge \bar\partial g + g^{-1}\bar\partial^2g = 0$, since a $1$-form wedged with itself vanishes and $\bar\partial^2 = 0$
(Proposition~\ref{complex-manifolds:proposition-dbar}).

(1) Two holomorphic frames differ by a holomorphic $g$, with $\bar\partial g = 0$, so by the computation in (2) the operators defined by the two frames agree; the Leibniz rule is clear, and $a_e = 0$ for a holomorphic frame, so
$\bar\partial_L$ is integrable. A section $fe$ is holomorphic if and only if $f$ is, that is, $\bar\partial f = 0$.

(3) Let $x \in X$, choose a coordinate polydisc $\Delta$ around $x$ on which $L$ has a smooth frame $e$, and write $a_e = \sum_jf_j\,d\bar z_j$. The condition $\bar\partial a_e = 0$ says
$\partial f_j/\partial\bar z_l = \partial f_l/\partial\bar z_j$, so by the $\bar\partial$-Poincar\'e lemma (Theorem~\ref{several-complex-variables:theorem-dbar-poincare}) there is a smooth $u$ on a smaller polydisc $\Delta'$
with $\bar\partial u = a_e$. The frame $e' = e^{-u}e$ satisfies $De' = (-\bar\partial u + a_e)e' = 0$. Cover $X$ by open sets $U_i$ with frames $e_i$ such that $De_i = 0$. On overlaps $e_i = g_{ij}e_j$ with $g_{ij}$ smooth and
nowhere zero, and $0 = De_i = \bar\partial g_{ij} \otimes e_j$, so the $g_{ij}$ are holomorphic and define a holomorphic structure on $L$ for which the $e_i$ are holomorphic frames. For a section $s = fe_i$ we
get $Ds = \bar\partial f \otimes e_i = \bar\partial_Ls$, so $D = \bar\partial_L$. For uniqueness, the holomorphic frames of any structure with $\bar\partial_L = D$ are the local frames killed by $D$, so the holomorphic
atlas is determined by $D$.
\end{proof}

\begin{example}\label{complex-manifolds:example-dbar-trivial-bundle}
On the smooth trivial line bundle $X \times \CC$, the $\bar\partial$-operators are the $D_a = \bar\partial + a$ with $a$ a $(0,1)$-form, integrable exactly when $\bar\partial a = 0$. A holomorphic isomorphism between the
structures defined by $D_a$ and $D_b$ is multiplication by a smooth nowhere vanishing function $g$ taking $D_a$-holomorphic sections to $D_b$-holomorphic ones; since
$D_b(gs) = g(D_as + (g^{-1}\bar\partial g + b - a)s)$, this happens if and only if $a - b = g^{-1}\bar\partial g$. In particular, if $a - b = \bar\partial u$ for a smooth function $u$, then $g = e^u$ is an
isomorphism. So $a \mapsto (X \times \CC, D_a)$ defines a map from the $\bar\partial$-closed $(0,1)$-forms modulo $\bar\partial$ of functions to $\Pic(X)$, whose image consists of line bundles that are trivial as smooth
bundles; compare the map $H^1(X, \mathcal{O}_X) \to \Pic(X)$ of Theorem~\ref{complex-manifolds:theorem-exponential}, and Section~\ref{complex-manifolds:section-dolbeault}, where the quotient above is the Dolbeault group
$H^{0,1}_{\bar\partial}(X)$.
\end{example}

\section{The holomorphic tangent and cotangent bundles}\label{complex-manifolds:section-tangent}

A complex manifold $X$ of dimension $n$ has three ``tangent bundles'': the real tangent bundle $TX$ of rank $2n$, and the two summands $T^{1,0}X$ and $T^{0,1}X$ of its complexification
$TX \otimes \CC = T^{1,0}X \oplus T^{0,1}X$ (Definition~\ref{complex-manifolds:definition-almost-complex}). This section explains their special properties. Only $T^{1,0}X$ is a \emph{holomorphic} bundle, because changes of holomorphic
coordinates act on it by holomorphic Jacobian matrices. At the same time it is nothing new: as a complex vector bundle it is $TX$, with $J$ as multiplication by $i$. Dually, the holomorphic cotangent bundle
$\Lambda^{1,0}X$ has as holomorphic sections exactly the holomorphic $1$-forms.

Recall that for a smooth complex-valued function $f$ in holomorphic coordinates $z^j = x^j + iy^j$,
\[
\frac{\partial}{\partial z^j} = \frac12\Bigl(\frac{\partial}{\partial x^j} - i\frac{\partial}{\partial y^j}\Bigr), \qquad \frac{\partial}{\partial\bar z^j} = \frac12\Bigl(\frac{\partial}{\partial x^j} + i\frac{\partial}{\partial y^j}\Bigr), \qquad
df = \sum_j\frac{\partial f}{\partial z^j}\,dz^j + \sum_j\frac{\partial f}{\partial\bar z^j}\,d\bar z^j,
\]
the last because both sides equal $\sum_j(\partial_{x^j}f\,dx^j + \partial_{y^j}f\,dy^j)$ after substituting $dz^j = dx^j + i\,dy^j$ and $d\bar z^j = dx^j - i\,dy^j$. Moreover $dz^k(\partial/\partial z^j) = \delta_{jk}$ and
$d\bar z^k(\partial/\partial z^j) = 0$; for instance $dz(\partial_z) = \frac12(dx + i\,dy)(\partial_x - i\partial_y) = \frac12(1 + 1) = 1$ and $d\bar z(\partial_z) = \frac12(1 - 1) = 0$. So $(dz^j, d\bar z^j)$ is the
basis of $T^*_pX \otimes \CC$ dual to the basis $(\partial/\partial z^j, \partial/\partial\bar z^j)$ of $T_pX \otimes \CC$.

\begin{proposition}\label{complex-manifolds:proposition-holomorphic-tangent}
Let $X$ be a complex manifold of dimension $n$, with its almost complex structure $J$, and let $z = (z^1, \ldots, z^n)$ and $w = (w^1, \ldots, w^n)$ be holomorphic coordinates on two overlapping charts.
\begin{enumerate}
\item On the overlap,
\[
\frac{\partial}{\partial z^j} = \sum_{k=1}^n\frac{\partial w^k}{\partial z^j}\,\frac{\partial}{\partial w^k}, \qquad dw^k = \sum_{j=1}^n\frac{\partial w^k}{\partial z^j}\,dz^j .
\]
Hence $T^{1,0}X$, with the local frames $(\partial/\partial z^j)$, is a holomorphic vector bundle of rank $n$, whose transition functions are the holomorphic Jacobian matrices $(\partial w^k/\partial z^j)$.
\item Likewise $\Lambda^{1,0}X$, with the local frames $(dz^j)$, is a holomorphic vector bundle, dual to $T^{1,0}X$, and its holomorphic sections are the holomorphic $1$-forms: as sheaves, $\Lambda^{1,0}X$ gives
$\Omega^1_X$. For a holomorphic function $f$, $df = \partial f = \sum_j(\partial f/\partial z^j)\,dz^j$ is such a section.
\item Regard $TX$ as a complex vector bundle by letting $i$ act as $J$. The map
\[
\phi \colon TX \to T^{1,0}X, \qquad \phi(v) = \tfrac12(v - iJv),
\]
is an isomorphism of complex vector bundles, with inverse $u \mapsto 2\operatorname{Re}u = u + \bar u$. In coordinates $\phi(\partial/\partial x^j) = \partial/\partial z^j$ and $\phi(\partial/\partial y^j) = i\,\partial/\partial z^j$.
\item $T^{0,1}X = \overline{T^{1,0}X}$ has the conjugate Jacobian matrices $(\overline{\partial w^k/\partial z^j})$ as transition functions with respect to the frames $(\partial/\partial\bar z^j)$, and $v \mapsto \frac12(v + iJv)$ is an
isomorphism of complex vector bundles from $(TX, -J)$ to $T^{0,1}X$.
\end{enumerate}
\end{proposition}

\begin{proof}
(1) Since the $w^k$ are holomorphic functions of $z$, $\partial w^k/\partial\bar z^j = 0$ by the Cauchy--Riemann equations (Proposition~\ref{complex-manifolds:proposition-dbar}), so the formula for $df$ above gives
$dw^k = \sum_j(\partial w^k/\partial z^j)\,dz^j$, and by conjugation $d\bar w^k = \sum_j\overline{(\partial w^k/\partial z^j)}\,d\bar z^j$. Expand the vector $\partial/\partial z^j$ in the basis
$(\partial/\partial w^k, \partial/\partial\bar w^k)$; the coefficients are the values of the dual basis on it. Now $dw^k(\partial/\partial z^j) = \partial w^k/\partial z^j$, while $d\bar w^k(\partial/\partial z^j) = 0$, since
$d\bar w^k$ is a combination of the $d\bar z^l$. This gives the formula for $\partial/\partial z^j$. The matrix $g = (\partial w^k/\partial z^j)$ has holomorphic entries, as derivatives of holomorphic functions are
holomorphic, and it is invertible, with inverse the Jacobian of $z$ as a function of $w$, by the chain rule. Two frames of the same bundle related by such matrices define a holomorphic structure: the transition
functions are holomorphic maps into $\GL_n(\CC)$, and the cocycle condition holds automatically, because they compare frames of one bundle (Definition~\ref{complex-manifolds:definition-holomorphic-bundle}).

(2) The frames $(dz^j)$ transform by the transpose of the same holomorphic matrix, so $\Lambda^{1,0}X$ is holomorphic, and $dz^k(\partial/\partial z^j) = \delta_{jk}$ identifies it with the dual of $T^{1,0}X$. A
section $\sum_jf_j\,dz^j$ is holomorphic exactly when the $f_j$ are holomorphic in every chart, which is the definition of a holomorphic $1$-form (Proposition~\ref{complex-manifolds:proposition-dbar}). For holomorphic
$f$ the $d\bar z^j$-terms of $df$ vanish, and the coefficients $\partial f/\partial z^j$ are holomorphic.

(3) $J$ acts on $TX \otimes \CC$ complex-linearly, and $J(v - iJv) = Jv + iv = i(v - iJv)$, so $\phi(v)$ lies in the $+i$-eigenspace $T^{1,0}X$. It is complex linear: $\phi(Jv) = \frac12(Jv + iv) = i\phi(v)$.
It is injective, since the real part of $\frac12(v - iJv)$ in $TX \otimes \CC = TX \oplus i\,TX$ is $\frac12v$. Both bundles have real rank $2n$, so $\phi$ is an isomorphism on each fibre, and it is smooth. For
$u = \phi(v)$, $u + \bar u = \frac12(v - iJv) + \frac12(v + iJv) = v$. In coordinates, $J\partial_{x^j} = \partial_{y^j}$ and $J\partial_{y^j} = -\partial_{x^j}$, so
$\phi(\partial_{x^j}) = \frac12(\partial_{x^j} - i\partial_{y^j}) = \partial/\partial z^j$ and $\phi(\partial_{y^j}) = \frac12(\partial_{y^j} + i\partial_{x^j}) = i\,\partial/\partial z^j$.

(4) Conjugate the formula of (1). For the last statement, $v + iJv$ lies in the $-i$-eigenspace, and $\frac12(-Jv + iv) = i \cdot \frac12(v + iJv)$, so the map is complex linear for the structure $-J$ on $TX$; it is
injective for the same reason as in (3).
\end{proof}

\begin{remark}\label{complex-manifolds:remark-holomorphic-tangent}
\begin{enumerate}
\item \emph{Why $T^{1,0}$ and not $T^{0,1}$.} The transition functions of $T^{0,1}X$ in (4) are antiholomorphic, so the frames $(\partial/\partial\bar z^j)$ do not make it a holomorphic bundle. Both bundles are
isomorphic to $TX$ as real bundles; the holomorphic structure is extra information that only $T^{1,0}X$ receives from the coordinates. This is why the holomorphic tangent bundle is written $T_X = T^{1,0}X$,
and its holomorphic sections are the \emph{holomorphic vector fields} $\sum_jf_j\,\partial/\partial z^j$, $f_j$ holomorphic.
\item \emph{Chern classes.} By (3), the Chern classes $c_k(X) = c_k(T^{1,0}X)$ are the Chern classes of $(TX, J)$, so they depend only on the almost complex structure. By (4) and Chapter~\ref{characteristic-classes},
$c_k(T^{0,1}X) = (-1)^kc_k(X)$. In particular $c_1(K_X) = c_1(\Lambda^n\Lambda^{1,0}X) = -c_1(X)$, since $\Lambda^{1,0}X$ is the dual of $T^{1,0}X$.
\item \emph{Real vector fields from holomorphic ones.} Under $\phi$, a holomorphic vector field $u = \sum_jf_j\,\partial/\partial z^j$ corresponds to the real vector field $V = 2\operatorname{Re}u$. Since $dz^j$ vanishes on
$T^{0,1}X$, $dz^j(V) = dz^j(u + \bar u) = f_j$. So in coordinates the integral curves of $V$ are the solutions of the complex differential equation $dz^j/dt = f_j(z)$. In dimension one, $f\,\partial_z$ with $f = a + ib$
corresponds to $a\,\partial_x + b\,\partial_y$.
\end{enumerate}
\end{remark}

\begin{example}\label{complex-manifolds:example-tangent-p1}
Let $X = \CC\PP^1 = \CC \cup \{\infty\}$, with the coordinate $z$ on $\CC$ and $w = 1/z$ near $\infty$. On the overlap $dw/dz = -1/z^2 = -w^2$, so by Proposition~\ref{complex-manifolds:proposition-holomorphic-tangent}(1)
\[
\frac{\partial}{\partial z} = -w^2\frac{\partial}{\partial w}, \qquad dz = -\frac{dw}{w^2}.
\]
\begin{enumerate}
\item \emph{Holomorphic vector fields.} A holomorphic vector field is $f(z)\,\partial_z$ on $\CC$, with $f$ entire, such that $f(z)\,\partial_z = -w^2f(1/w)\,\partial_w$ extends holomorphically over $w = 0$. If
$f = c_0 + c_1z + c_2z^2$, then $-w^2f(1/w) = -(c_0w^2 + c_1w + c_2)$, which extends. Conversely, if $g(w) = w^2f(1/w)$ extends holomorphically to $w = 0$, it is bounded near $0$, so $|f(z)| \leq C|z|^2$ for $|z|$
large; as $f$ is also bounded on compact sets, $|f(z)| \leq C^{\prime}(1 + |z|)^2$, and $f$ is a polynomial of degree at most $2$ by Liouville's theorem (Corollary~\ref{complex-analysis:corollary-liouville}(1)). So the
holomorphic vector fields form the $3$-dimensional space spanned by $\partial_z$, $z\,\partial_z$, $z^2\,\partial_z$. By Remark~\ref{complex-manifolds:remark-holomorphic-tangent}(3) their real flows solve
$\dot z = 1$, $\dot z = z$ and $\dot z = z^2$, namely
\[
z \mapsto z + t, \qquad z \mapsto e^tz, \qquad z \mapsto \frac{z}{1 - tz},
\]
(check the last: $\frac{d}{dt}\frac{z}{1 - tz} = \frac{z^2}{(1 - tz)^2}$). All three are one-parameter groups of M\"obius transformations, defined on all of $\CC\PP^1$ for all $t$. Replacing $f$ by $if$ gives for instance the rotations
$z \mapsto e^{it}z$, and composing these flows gives all the maps $z \mapsto z + b$, $z \mapsto az$ and $z \mapsto z/(1 - cz)$ with $b, c \in \CC$ and $a \in \CC^\times$. So the holomorphic vector fields are
infinitesimal automorphisms of $\CC\PP^1$.
\item \emph{No holomorphic $1$-forms.} A holomorphic $1$-form is $f(z)\,dz$ on $\CC$ such that $-f(1/w)w^{-2}\,dw$ extends over $w = 0$. Then $|f(z)| \leq C|z|^{-2}$ for $|z|$ large, so $f$ is bounded, hence constant
by Liouville's theorem, and the constant is $\lim_{z \to \infty}f(z) = 0$. So $H^0(\CC\PP^1, \Omega^1) = 0$. On a complex torus $\CC/\Lambda$, by contrast, $dz$ and $\partial_z$ are invariant under translations, so they
descend to nowhere vanishing sections: $T^{1,0}$ and $\Omega^1$ of a torus are trivial.
\item \emph{$T^{1,0}\CC\PP^1 \cong \mathcal{O}(2)$ and $K_{\CC\PP^1} \cong \mathcal{O}(-2)$.} With homogeneous coordinates $[x_0 : x_1]$, $z = x_1/x_0$ on $U_0$ and $w = x_0/x_1$ on $U_1$, a vector field $s_0\,\partial_z = s_1\,\partial_w$
has $s_1 = -(x_0/x_1)^2s_0$. The sections of $\mathcal{O}(2)$ satisfy $s_1 = (x_0/x_1)^2s_0$ (Exercise~\ref{complex-manifolds:exercise-sections-ok}). The two cocycles differ by the constant $-1$, which is
removed by changing the frame on $U_1$ from $\partial_w$ to $-\partial_w$; so $T^{1,0}\CC\PP^1 \cong \mathcal{O}(2)$, in accordance with (1): the homogeneous quadrics $x_0^2$, $x_0x_1$, $x_1^2$ form a basis of
the sections of $\mathcal{O}(2)$. Dually $K_{\CC\PP^1} = \Omega^1 \cong \mathcal{O}(-2)$, which has no sections, in accordance with (2); this is the case $n = 1$ of $K_{\CC\PP^n} \cong \mathcal{O}(-n-1)$
(Proposition~\ref{complex-manifolds:proposition-adjunction}).
\end{enumerate}
\end{example}

\begin{remark}\label{complex-manifolds:remark-holomorphic-symplectic}
The holomorphic cotangent bundle has, like the real one (Chapter~\ref{symplectic-geometry}), a canonical symplectic form, now holomorphic. Its total space $T^*_{\mathrm{hol}}X$ (the total space of $\Lambda^{1,0}X$) is a complex
manifold of dimension $2n$, with holomorphic coordinates $(z^1, \ldots, z^n, p_1, \ldots, p_n)$ for the covector $\sum_jp_j\,dz^j$ at the point $z$; in another chart the $p$'s change by the holomorphic matrix of
Proposition~\ref{complex-manifolds:proposition-holomorphic-tangent}(2). The holomorphic $1$-form $\theta = \sum_jp_j\,dz^j$ is independent of the chart: at the covector $\xi = \sum_jp_j\,dz^j$ it is $\xi \circ d\pi$, with $\pi$ the
projection to $X$, exactly as for the tautological form (Definition~\ref{symplectic-geometry:definition-tautological} and Proposition~\ref{symplectic-geometry:proposition-tautological}(1)), and $\omega = d\theta = \sum_jdp_j \wedge dz^j$ is a closed holomorphic $2$-form that is nondegenerate at every point. Complex manifolds with such a form are
called \emph{holomorphic symplectic}. Another source of examples: on a complex surface ($n = 2$) with trivial canonical bundle, a nowhere vanishing section of $K_X = \Omega^2_X$ is a holomorphic $2$-form that is
nondegenerate at every point, because a nonzero alternating form on a $2$-dimensional space is nondegenerate. So complex tori of dimension $2$ and K3 surfaces (Example~\ref{complex-manifolds:example-adjunction}) are
holomorphic symplectic.
\end{remark}

\section{Divisors}\label{complex-manifolds:section-divisors}

\begin{definition}\label{complex-manifolds:definition-divisor}
A \emph{divisor} on $X$ is a locally finite formal sum $D = \sum_ia_iZ_i$ with $a_i \in \ZZ$ and $Z_i$ irreducible analytic hypersurfaces; they form the group $\mathrm{Div}(X)$. A \emph{meromorphic function}
is a section of the sheaf $\mathcal{M}_X$ of quotients $f/g$ of holomorphic functions with $g \not\equiv 0$; the order $\operatorname{ord}_Z(f) \in \ZZ$ of a meromorphic function along an irreducible
hypersurface $Z$ is defined by writing $f = u\,h^m$ near a regular point of $Z$, where $h$ is a local defining function of $Z$ and $u$ is a unit, using that $\mathcal{O}_{X,x}$ is a unique factorisation domain
(Theorem~\ref{several-complex-variables:theorem-local-ring}). The divisor of $f$ is $(f) = \sum_Z\operatorname{ord}_Z(f)\,Z$, and divisors of this form are \emph{principal}.
\end{definition}

\begin{proposition}\label{complex-manifolds:proposition-divisor-line-bundle}
There is a homomorphism $\mathrm{Div}(X) \to \Pic(X)$, $D \mapsto \mathcal{O}(D)$, defined by taking local defining functions $f_i$ for $D$ on a cover and the cocycle $g_{ij} = f_i/f_j \in \mathcal{O}^*$. Its kernel
is the group of principal divisors, and $\mathcal{O}(D)$ has a canonical meromorphic section with divisor $D$; for $D$ effective (all $a_i \geq 0$) this section is holomorphic and vanishes exactly on $D$.
Conversely a nonzero meromorphic section of a line bundle $L$ has a divisor $D$ with $L \cong \mathcal{O}(D)$.
\end{proposition}

\begin{proof}
The functions $f_i/f_j$ are nowhere vanishing and holomorphic on overlaps, because two local defining functions of the same divisor differ by a unit
(Theorem~\ref{several-complex-variables:theorem-local-ring}), and they clearly form a cocycle; changing the $f_i$ changes the cocycle by a coboundary. The local functions $f_i$ glue to a meromorphic section
of $\mathcal{O}(D)$, whose divisor is $D$. If $\mathcal{O}(D)$ is trivial, the cocycle is a coboundary $f_i/f_j = u_i/u_j$, so the $f_i/u_i$ glue to a global meromorphic function with divisor $D$; conversely,
principal divisors give trivial cocycles. The last statement follows by taking the divisor of the section and comparing cocycles.
\end{proof}

\begin{example}\label{complex-manifolds:example-divisors-projective}
On $\CC\PP^n$ every irreducible hypersurface is the zero set of an irreducible homogeneous polynomial, and $\mathcal{O}(D) \cong \mathcal{O}(\deg D)$; the map $\deg \colon \mathrm{Div}(\CC\PP^n) \to \ZZ$
induces the isomorphism $\Pic(\CC\PP^n) \cong \ZZ$, which also follows from Theorem~\ref{complex-manifolds:theorem-exponential} since $H^1(\CC\PP^n, \mathcal{O}) = H^2(\CC\PP^n, \mathcal{O}) = 0$
(Section~\ref{complex-manifolds:section-dolbeault}). On a complex torus $X = \CC^n/\Lambda$, by contrast, $H^1(X, \mathcal{O}_X) \neq 0$ and $\Pic^0(X)$ is a nontrivial complex torus, the dual torus.
\end{example}

\begin{counterexample}\label{complex-manifolds:counterexample-topologically-trivial}
Holomorphic line bundles are not determined by their topology. Let $X = \CC/\Lambda$ with $\Lambda = \ZZ\omega_1 + \ZZ\omega_2$ a lattice, numbered so that $\operatorname{Im}(\omega_2/\omega_1) > 0$, $\pi \colon \CC \to X$ the projection, and $p \neq q$ points of $X$. The line bundle
$L = \mathcal{O}(p - q) = \mathcal{O}(p) \otimes \mathcal{O}(q)^\vee$ is trivial as a smooth bundle but not as a holomorphic bundle.

\emph{Smooth triviality.} Lift $q - p$ to $c \in \CC$ and let $\tau(x) = x + c$, a biholomorphism of $X$ with $\tau(p) = q$. A local defining function of $q$ composed with $\tau$ is a local defining function of $p$, so
$\tau^*\mathcal{O}(q) \cong \mathcal{O}(p)$. The translations $x \mapsto x + tc$, $0 \leq t \leq 1$, form a homotopy from the identity to $\tau$, so by Theorem~\ref{vector-bundles:theorem-line-bundles} and homotopy
invariance (Theorem~\ref{singular-homology:theorem-axioms}), $c_1(\mathcal{O}(p)) = \tau^*c_1(\mathcal{O}(q)) = c_1(\mathcal{O}(q))$. Hence $c_1(L) = 0$ and $L$ is smoothly trivial, by the same theorem.

\emph{Holomorphic nontriviality.} If $L$ were holomorphically trivial, the divisor $p - q$ would be principal (Proposition~\ref{complex-manifolds:proposition-divisor-line-bundle}): there would be a meromorphic
function $f$ on $X$ whose only pole is a simple pole at $q$. Then $F = f \circ \pi$ is meromorphic on $\CC$, invariant under translation by $\Lambda$, with simple poles exactly at the points of $\pi^{-1}(q)$. Choose a lift
$\tilde q$ of $q$ and let $P$ be the parallelogram with vertices $a$, $a + \omega_1$, $a + \omega_1 + \omega_2$, $a + \omega_2$, where $a = \tilde q - \frac12(\omega_1 + \omega_2)$. A boundary point
$\tilde q + (s - \frac12)\omega_1 + (t - \frac12)\omega_2$ has $s$ or $t$ in $\{0, 1\}$, so it is not in $\tilde q + \Lambda$; and $\tilde q$ is the only point of $\tilde q + \Lambda$ inside $P$. The boundary $\partial P$, traversed through the vertices in the order listed, is positively oriented, with winding number $1$ about
$\tilde q$ and $0$ about the points outside $P$. By the residue theorem
(Theorem~\ref{complex-analysis:theorem-residues}), $\int_{\partial P}F\,dz = 2\pi i\operatorname{Res}_{\tilde q}F \neq 0$, the residue at a simple pole being nonzero. On the other hand the integrals over opposite sides of $P$
cancel: the side from $a + \omega_2$ to $a + \omega_1 + \omega_2$ is the translate by $\omega_2$ of the side from $a$ to $a + \omega_1$, traversed in the opposite direction, and $F(z + \omega_2) = F(z)$; similarly for the
other pair. So $\int_{\partial P}F\,dz = 0$, a contradiction.

Thus $L$ is a nontrivial element of $\Pic^0(X)$, the kernel of $c_1$ in Theorem~\ref{complex-manifolds:theorem-exponential}. The same argument shows that no meromorphic function on $X$ has a single pole which is simple.
\end{counterexample}

\section{Hermitian metrics and the Chern connection}\label{complex-manifolds:section-chern-connection}

\begin{theorem}\label{complex-manifolds:theorem-chern-connection}
Let $E$ be a holomorphic vector bundle with a Hermitian metric $h$. There is a unique connection $\nabla$ on $E$, the \emph{Chern connection}, that is metric and satisfies $\nabla^{0,1} = \bar\partial$. In a
holomorphic frame with $H_{ij} = h(e_i, e_j)$ one has $\theta = \overline{H}^{-1}\partial\overline{H}$ and curvature $F = \bar\partial\theta$, a $(1,1)$-form with values in $\End E$. For a line bundle with local
weight $h = e^{-\varphi}$ in a holomorphic frame, $F = \partial\bar\partial\varphi$ and
\[
c_1(L) = \left[\frac{i}{2\pi}F\right] = \left[\frac{i}{2\pi}\partial\bar\partial\varphi\right] \in H^2(X; \RR),
\]
a class of type $(1,1)$.
\end{theorem}

\begin{proof}
Write $\nabla = \nabla^{1,0} + \nabla^{0,1}$. The condition $\nabla^{0,1} = \bar\partial$ fixes the $(0,1)$-part, and the metric condition
$d\,h(s, t) = h(\nabla s, t) + h(s, \nabla t)$ splits into its $(1,0)$- and $(0,1)$-parts, the first of which reads $\partial h(s, t) = h(\nabla^{1,0}s, t) + h(s, \bar\partial t)$ and determines $\nabla^{1,0}$; this
proves uniqueness, and the resulting formula in a holomorphic frame, $\theta = \overline{H}^{-1}\partial\overline{H}$, defines a connection with the two properties, which proves existence. Then
$F = d\theta + \theta \wedge \theta$ has no $(0,2)$-part because $\bar\partial^2 = 0$ and $\nabla^{0,1} = \bar\partial$, and no $(2,0)$-part by the metric condition; the formula $F = \bar\partial\theta$ follows. For a
line bundle $H = e^{-\varphi}$ is real, so $\theta = H^{-1}\partial H = \partial\log H = -\partial\varphi$ and $F = \bar\partial\theta = -\bar\partial\partial\varphi = \partial\bar\partial\varphi$; the Chern class is
computed by Chern--Weil theory (Proposition~\ref{characteristic-classes:proposition-chern-properties}).
\end{proof}

\begin{definition}\label{complex-manifolds:definition-positive}
A Hermitian line bundle $(L, h)$ is \emph{positive} if the real $(1,1)$-form $\frac{i}{2\pi}F$ is positive definite, that is, if it is the fundamental form of a Hermitian metric on $X$
(Chapter~\ref{kahler}); $L$ is positive if it admits such a metric. A line bundle is \emph{ample} if a power of it gives an embedding of $X$ in projective space
(Chapter~\ref{kodaira}).
\end{definition}

\begin{example}\label{complex-manifolds:example-fubini-study}
On $\CC\PP^n$, give $\mathcal{O}(-1)$ the metric induced from the standard Hermitian metric of $\CC^{n+1}$. On $U_0 = \{x_0 \neq 0\}$, with coordinates $w_j = x_j/x_0$, the section
$\sigma(w) = (1, w_1, \ldots, w_n)$ is a holomorphic frame with $h(\sigma, \sigma) = 1 + |w|^2 = e^{-\varphi}$, so by Theorem~\ref{complex-manifolds:theorem-chern-connection}
$\frac{i}{2\pi}F = -\frac{i}{2\pi}\partial\bar\partial\log(1 + |w|^2) = -\omega_{FS}$, where $\omega_{FS}$, the \emph{Fubini--Study form}, is the global $(1,1)$-form given on $U_0$ by
$\frac{i}{2\pi}\partial\bar\partial\log(1 + |w|^2)$ (on $U_j$ one uses $\log(|x|^2/|x_j|^2)$ in place of $\log(1 + |w|^2) = \log(|x|^2/|x_0|^2)$; on overlaps these differ by $\log|x_k/x_j|^2$, the logarithm of the squared modulus of a nowhere
vanishing holomorphic function, which $\partial\bar\partial$ kills); hence
$\mathcal{O}(1)$ is positive, with $\frac{i}{2\pi}F = \omega_{FS}$ and $c_1(\mathcal{O}(1)) = [\omega_{FS}]$, the positive generator of $H^2(\CC\PP^n; \ZZ)$
(Example~\ref{characteristic-classes:example-tautological}).
\end{example}

\section{The canonical bundle and adjunction}\label{complex-manifolds:section-adjunction}

\begin{proposition}[Adjunction]\label{complex-manifolds:proposition-adjunction}
Let $Z \subseteq X$ be a smooth hypersurface with associated line bundle $\mathcal{O}(Z)$. Then the normal bundle is $N_{Z/X} \cong \mathcal{O}(Z)|_Z$, and
\[
K_Z \cong \left(K_X \otimes \mathcal{O}(Z)\right)\big|_Z .
\]
For a smooth hypersurface of degree $d$ in $\CC\PP^n$ one gets $K_Z \cong \mathcal{O}(d - n - 1)|_Z$.
\end{proposition}

\begin{proof}
A local defining function $f_i$ of $Z$ gives, by $df_i$, an isomorphism of the conormal bundle $N^\vee_{Z/X}$ with $\mathcal{O}(-Z)|_Z$: the transition functions of the conormal bundle are the
$f_i/f_j$ restricted to $Z$, computed from $df_i = (f_i/f_j)\,df_j$ modulo $T^*Z$. Taking determinants in the exact sequence
$0 \to N^\vee_{Z/X} \to \Omega^1_X|_Z \to \Omega^1_Z \to 0$ gives $K_X|_Z \cong K_Z \otimes N^\vee_{Z/X} \cong K_Z \otimes \mathcal{O}(-Z)|_Z$, which is the formula. For $\CC\PP^n$ use
$K_{\CC\PP^n} \cong \mathcal{O}(-n-1)$, which follows from the Euler sequence (Example~\ref{vector-bundles:example-euler-sequence}) by taking determinants, and $\mathcal{O}(Z) = \mathcal{O}(d)$.
\end{proof}

\begin{example}\label{complex-manifolds:example-adjunction}
A smooth plane curve of degree $d$ in $\CC\PP^2$ has $K_Z \cong \mathcal{O}(d - 3)|_Z$, whence genus $g = (d-1)(d-2)/2$ by Riemann--Roch (Chapter~\ref{curves}); a smooth quadric surface in $\CC\PP^3$ has
$K \cong \mathcal{O}(-2)$ and is $\CC\PP^1 \times \CC\PP^1$; a smooth quartic surface in $\CC\PP^3$ has trivial canonical bundle and is a K3 surface, one of the basic examples in Hodge theory.
\end{example}

\section{Dolbeault Cohomology}\label{complex-manifolds:section-dolbeault}
Dolbeault cohomology is the cohomology of the $\bar\partial$-operator on forms of type $(p,q)$. By the Dolbeault theorem
(Theorem~\ref{complex-manifolds:theorem-dolbeault}) it computes the sheaf cohomology of the holomorphic $p$-forms, and on a compact complex
manifold harmonic theory for $\bar\partial$ makes it finite-dimensional and gives Serre duality
(Theorem~\ref{complex-manifolds:theorem-serre-duality}). The numbers $h^{p,q} = \dim H^q(X, \Omega^p_X)$ are the \emph{Hodge numbers}; for
compact K\"ahler manifolds they are the dimensions of the pieces of the Hodge decomposition (Chapter~\ref{kahler}).
References: \cite[Chapter 0]{griffiths-harris}, \cite[Chapter 4]{voisin-hodge-1}, \cite[Chapters 2--4]{huybrechts-complex},
\cite{wells-differential-analysis}.

\noindent\textbf{Prerequisites.} Section~\ref{complex-manifolds:section-holomorphic-bundles} (Holomorphic Vector Bundles) and Chapter~\ref{elliptic-operators}
(Elliptic Operators on Compact Manifolds).

$X$ is a complex manifold of dimension $n$ and $E$ a holomorphic vector bundle on it.

\section{The Dolbeault theorem}\label{complex-manifolds:section-dolbeault-theorem}

\begin{definition}\label{complex-manifolds:definition-dolbeault}
Let $\mathcal{A}^{p,q}(E)$ be the sheaf of smooth $E$-valued $(p,q)$-forms and $\mathcal{A}^{p,q}(X, E)$ its global sections. In a holomorphic frame of $E$ set $\bar\partial_E(\sum_j\alpha_j e_j) = \sum_j\bar\partial\alpha_j\,e_j$;
this is independent of the frame because the transition functions are holomorphic, and $\bar\partial_E^2 = 0$. The \emph{Dolbeault cohomology} is
$H^{p,q}_{\bar\partial}(X, E) = H^q(\mathcal{A}^{p,\bullet}(X, E), \bar\partial_E)$; for $E$ trivial we write $H^{p,q}_{\bar\partial}(X)$.
\end{definition}

\begin{lemma}\label{complex-manifolds:lemma-dbar-poincare-pq}
Let $\Delta' \subseteq \Delta$ be polydiscs with $\overline{\Delta'} \subseteq \Delta$ and $\alpha$ a smooth $(p,q)$-form on $\Delta$ with $q \geq 1$ and $\bar\partial\alpha = 0$. Then $\alpha = \bar\partial\beta$ on $\Delta'$ for some
smooth $(p, q-1)$-form $\beta$.
\end{lemma}

\begin{proof}
Writing $\alpha = \sum_I dz^I \wedge \alpha_I$ with $\alpha_I$ of type $(0,q)$, the condition $\bar\partial\alpha = 0$ holds componentwise, since $\bar\partial(dz^I) = 0$; so we may assume $p = 0$. For $q = 1$ this is
Theorem~\ref{several-complex-variables:theorem-dbar-poincare}. For general $q$ the same induction applies: if $\alpha$ involves only $d\bar z^1, \ldots, d\bar z^k$, write
$\alpha = d\bar z^k \wedge \alpha_1 + \alpha_2$ with $\alpha_1, \alpha_2$ not involving $d\bar z^k, \ldots, d\bar z^n$; the equation $\bar\partial\alpha = 0$ implies that the coefficients of $\alpha_1$ are holomorphic in
$z_{k+1}, \ldots, z_n$. Solving $\partial g/\partial\bar z_k = $ (coefficient of $\alpha_1$) with the Cauchy transform of Lemma~\ref{several-complex-variables:lemma-cauchy-transform}, after a cut-off in $z_k$
alone, gives a form $\gamma$ with $\alpha - \bar\partial\gamma$ involving only $d\bar z^1, \ldots, d\bar z^{k-1}$ on a slightly smaller polydisc; by induction on $k$ we conclude. The details are exactly as in the proof of
Theorem~\ref{several-complex-variables:theorem-dbar-poincare}; see \cite[Chapter 0, Section 2]{griffiths-harris}.
\end{proof}

\begin{theorem}[Dolbeault]\label{complex-manifolds:theorem-dolbeault}
The sequence of sheaves $0 \to \Omega^p_X(E) \to \mathcal{A}^{p,0}(E) \xrightarrow{\bar\partial_E} \mathcal{A}^{p,1}(E) \xrightarrow{\bar\partial_E} \cdots$ is a fine resolution of the sheaf
$\Omega^p_X(E) = \Omega^p_X \otimes E$ of holomorphic $E$-valued $p$-forms. Consequently
\[
H^{p,q}_{\bar\partial}(X, E) \cong H^q(X, \Omega^p_X \otimes E).
\]
\end{theorem}

\begin{proof}
Exactness at $\mathcal{A}^{p,0}$ says that the $E$-valued $(p,0)$-forms killed by $\bar\partial_E$ are the holomorphic ones (Proposition~\ref{complex-manifolds:proposition-dbar}), and exactness at
$\mathcal{A}^{p,q}$ for $q \geq 1$ holds on stalks by Lemma~\ref{complex-manifolds:lemma-dbar-poincare-pq}, applied in a holomorphic frame of $E$ on small polydiscs. The sheaves $\mathcal{A}^{p,q}(E)$ are modules over the
sheaf of smooth functions, hence fine and acyclic (Theorem~\ref{sheaves:theorem-soft-acyclic}); apply Proposition~\ref{sheaves:proposition-hypercohomology}(3).
\end{proof}

\begin{theorem}[Holomorphic Poincar\'e lemma]\label{complex-manifolds:theorem-holomorphic-de-rham}
The holomorphic de Rham complex $0 \to \CC_X \to \Omega^0_X \xrightarrow{\partial} \Omega^1_X \xrightarrow{\partial} \cdots \to \Omega^n_X \to 0$ is exact. Hence
$H^k(X; \CC) \cong \mathbb{H}^k(X, \Omega^\bullet_X)$, and there is a spectral sequence (the \emph{Fr\"olicher} or \emph{Hodge-to-de Rham} spectral sequence)
\[
E_1^{p,q} = H^q(X, \Omega^p_X) \Rightarrow H^{p+q}(X; \CC).
\]
\end{theorem}

\begin{proof}
On holomorphic forms $d = \partial$. Consider the double complex of sheaves $\mathcal{A}^{p,q}$ with differentials $\partial$ and $\bar\partial$, whose total complex is the complexified de Rham complex. By
Theorem~\ref{complex-manifolds:theorem-dolbeault} each column $(\mathcal{A}^{p,\bullet}, \bar\partial)$ is a resolution of $\Omega^p_X$; by Lemma~\ref{homological-algebra:lemma-acyclic-assembly}, applied to stalks, the
inclusion $\Omega^\bullet_X \to \mathrm{Tot}(\mathcal{A}^{\bullet,\bullet}) = \mathcal{A}^\bullet_\CC$ is a quasi-isomorphism. The de Rham complex is a resolution of $\CC_X$
(Theorem~\ref{de-rham:theorem-de-rham}), so $\Omega^\bullet_X$ has cohomology sheaf $\CC_X$ in degree $0$ and none elsewhere, which is the exactness. The isomorphism with $H^k(X; \CC)$ follows from
Proposition~\ref{sheaves:proposition-hypercohomology}(1) and Theorem~\ref{de-rham:theorem-de-rham}, and the spectral sequence is the first one of
Proposition~\ref{sheaves:proposition-hypercohomology}(2).
\end{proof}

\begin{proposition}\label{complex-manifolds:proposition-functoriality}
Let $f \colon X \to Y$ be a holomorphic map of complex manifolds. Pullback of forms maps $\mathcal{A}^{p,q}(Y)$ into $\mathcal{A}^{p,q}(X)$ and commutes with $\partial$ and $\bar\partial$, so it induces
maps $f^* \colon H^{p,q}_{\bar\partial}(Y) \to H^{p,q}_{\bar\partial}(X)$, with $(g \circ f)^* = f^* \circ g^*$. The wedge product induces products
$H^{p,q}_{\bar\partial}(X) \times H^{p',q'}_{\bar\partial}(X) \to H^{p+p',q+q'}_{\bar\partial}(X)$, making $\bigoplus_{p,q}H^{p,q}_{\bar\partial}(X)$ a graded-commutative algebra (for the total degree $p + q$), and
$f^*$ is a homomorphism of algebras.
\end{proposition}

\begin{proof}
In local holomorphic coordinates $w_j$ on $Y$, $f^*dw_j = d(w_j \circ f) = \partial(w_j \circ f)$ because $w_j \circ f$ is holomorphic, so $f^*dw_j$ has type $(1,0)$, and by conjugation $f^*d\bar w_j$ has type $(0,1)$;
hence $f^*$ preserves types. Since $f^*d = df^*$ and $d = \partial + \bar\partial$ with $\partial$ of type $(1,0)$ and $\bar\partial$ of type $(0,1)$ (Proposition~\ref{complex-manifolds:proposition-dbar}), comparing the
components of type $(p+1,q)$ and $(p,q+1)$ of $f^*d\alpha = df^*\alpha$ for $\alpha$ of type $(p,q)$ gives $f^*\partial\alpha = \partial f^*\alpha$ and $f^*\bar\partial\alpha = \bar\partial f^*\alpha$. So $f^*$ maps
$\bar\partial$-closed forms to $\bar\partial$-closed forms and $\bar\partial$-exact forms to $\bar\partial$-exact forms. Functoriality holds already for forms.

Similarly, taking the component of type $(p+p',q+q'+1)$ of the Leibniz rule $d(\alpha \wedge \beta) = d\alpha \wedge \beta + (-1)^{p+q}\alpha \wedge d\beta$ gives
$\bar\partial(\alpha \wedge \beta) = \bar\partial\alpha \wedge \beta + (-1)^{p+q}\alpha \wedge \bar\partial\beta$. Hence the wedge product of $\bar\partial$-closed forms is $\bar\partial$-closed, and if $\bar\partial\beta = 0$
then $\bar\partial\gamma \wedge \beta = \bar\partial(\gamma \wedge \beta)$ is exact; so the product is well defined on classes. Graded commutativity and compatibility with $f^*$ hold for forms.
\end{proof}

\begin{definition}\label{complex-manifolds:definition-bott-chern}
The \emph{Bott--Chern} and \emph{Aeppli cohomology} groups of $X$ are
\[
H^{p,q}_{BC}(X) = \frac{\{\alpha \in \mathcal{A}^{p,q}(X) : \partial\alpha = \bar\partial\alpha = 0\}}{\partial\bar\partial\mathcal{A}^{p-1,q-1}(X)}, \qquad
H^{p,q}_{A}(X) = \frac{\{\alpha \in \mathcal{A}^{p,q}(X) : \partial\bar\partial\alpha = 0\}}{\partial\mathcal{A}^{p-1,q}(X) + \bar\partial\mathcal{A}^{p,q-1}(X)}.
\]
\end{definition}

\begin{lemma}\label{complex-manifolds:lemma-bott-chern-maps}
The identity on forms induces natural linear maps
\[
H^{p,q}_{BC}(X) \to H^{p,q}_{\bar\partial}(X), \qquad H^{p,q}_{BC}(X) \to H^{p+q}_{\mathrm{dR}}(X; \CC), \qquad H^{p,q}_{\bar\partial}(X) \to H^{p,q}_{A}(X).
\]
Complex conjugation induces isomorphisms $H^{p,q}_{BC}(X) \cong H^{q,p}_{BC}(X)$ of real vector spaces, conjugate-linear over $\CC$.
\end{lemma}

\begin{proof}
The quotients in Definition~\ref{complex-manifolds:definition-bott-chern} make sense because $\partial\bar\partial\beta$ is $\partial$- and $\bar\partial$-closed and $\partial\alpha = 0$, $\bar\partial\alpha = 0$ imply
$\partial\bar\partial\alpha = 0$ (using $\partial^2 = \bar\partial^2 = 0$ and $\partial\bar\partial = -\bar\partial\partial$). If $\partial\alpha = \bar\partial\alpha = 0$, then $\alpha$ is $\bar\partial$-closed and
$d\alpha = 0$; and $\partial\bar\partial\beta = -\bar\partial(\partial\beta)$ is $\bar\partial$-exact, while $\partial\bar\partial\beta = d(\bar\partial\beta)$ is $d$-exact because $\bar\partial^2\beta = 0$. This gives the first
two maps. For the third, a $\bar\partial$-closed form is $\partial\bar\partial$-closed, and $\bar\partial$-exact forms lie in the denominator of $H_A$. Finally $\overline{\partial\alpha} = \bar\partial\bar\alpha$, so conjugation
preserves the conditions defining $H_{BC}$ and maps $\partial\bar\partial\mathcal{A}^{p-1,q-1}$ onto $\bar\partial\partial\mathcal{A}^{q-1,p-1} = \partial\bar\partial\mathcal{A}^{q-1,p-1}$.
\end{proof}

\begin{remark}\label{complex-manifolds:remark-bott-chern}
Unlike $H^{p,q}_{\bar\partial}$, the Bott--Chern groups have a definition symmetric in $p$ and $q$, and they map to de Rham cohomology. On a compact K\"ahler manifold the $\partial\bar\partial$-lemma
(Chapter~\ref{kahler}) says that a $d$-closed form of pure type which is $d$-exact or $\bar\partial$-exact is of the form $\partial\bar\partial\beta$; this is exactly the injectivity of the first two
maps of Lemma~\ref{complex-manifolds:lemma-bott-chern-maps}. On general compact complex manifolds these maps need not be injective, and Bott--Chern cohomology is a finer invariant than Dolbeault or de Rham
cohomology.
\end{remark}

\section{Harmonic theory for $\bar\partial$}\label{complex-manifolds:section-harmonic}

\begin{definition}\label{complex-manifolds:definition-dbar-laplacian}
Let $X$ be compact with a Hermitian metric and $E$ a Hermitian holomorphic bundle. With the $L^2$ inner products on $\mathcal{A}^{p,q}(X, E)$, let $\bar\partial_E^*$ be the formal adjoint of $\bar\partial_E$
(Definition~\ref{elliptic-operators:definition-formal-adjoint}) and $\Box_E = \bar\partial_E\bar\partial_E^* + \bar\partial_E^*\bar\partial_E$, the \emph{$\bar\partial$-Laplacian}. The space of
\emph{$\bar\partial$-harmonic forms} is $\mathcal{H}^{p,q}(X, E) = \ker\Box_E$.
\end{definition}

\begin{lemma}\label{complex-manifolds:lemma-harmonic-closed}
In the situation of Definition~\ref{complex-manifolds:definition-dbar-laplacian}, a form $\alpha \in \mathcal{A}^{p,q}(X, E)$ is $\Box_E$-harmonic if and only if $\bar\partial_E\alpha = 0$ and $\bar\partial_E^*\alpha = 0$. In
particular, for $q = 0$ the harmonic forms are exactly the holomorphic sections of $\Omega^p_X \otimes E$, for every choice of the metrics, and
\[
\mathcal{H}^{p,0}(X, E) = H^{p,0}_{\bar\partial}(X, E) = H^0(X, \Omega^p_X \otimes E).
\]
\end{lemma}

\begin{proof}
Since $X$ is compact, the defining property of formal adjoints gives $\langle\Box_E\alpha, \alpha\rangle = \|\bar\partial_E\alpha\|^2 + \|\bar\partial_E^*\alpha\|^2$ in $L^2$, so $\Box_E\alpha = 0$ forces
$\bar\partial_E\alpha = \bar\partial_E^*\alpha = 0$; the converse is clear from the definition of $\Box_E$. For $q = 0$, $\bar\partial_E^*$ maps $\mathcal{A}^{p,0}$ to $\mathcal{A}^{p,-1} = 0$, so the harmonic
$(p,0)$-forms are the $\bar\partial_E$-closed ones, which are the holomorphic sections of $\Omega^p_X \otimes E$ (Proposition~\ref{complex-manifolds:proposition-dbar}, in a holomorphic frame of $E$); and
$H^{p,0}_{\bar\partial}(X, E)$ is the kernel of $\bar\partial_E$ on $\mathcal{A}^{p,0}(X, E)$, as there are no forms of type $(p,-1)$.
\end{proof}

\begin{theorem}\label{complex-manifolds:theorem-dbar-hodge}
$\Box_E$ is elliptic, with principal symbol $\frac12|\xi|^2$. Consequently $\mathcal{H}^{p,q}(X, E)$ is finite-dimensional, there is an orthogonal decomposition
$\mathcal{A}^{p,q}(X, E) = \mathcal{H}^{p,q}(X, E) \oplus \bar\partial_E\mathcal{A}^{p,q-1}(X, E) \oplus \bar\partial_E^*\mathcal{A}^{p,q+1}(X, E)$, and every Dolbeault class has a unique harmonic
representative:
\[
H^q(X, \Omega^p_X \otimes E) \cong H^{p,q}_{\bar\partial}(X, E) \cong \mathcal{H}^{p,q}(X, E).
\]
In particular the Hodge numbers $h^{p,q}(X) = \dim H^q(X, \Omega^p_X)$ of a compact complex manifold are finite.
\end{theorem}

\begin{proof}
The symbol of $\bar\partial$ is $i\xi^{0,1} \wedge -$ and that of $\bar\partial^*$ is $-i\iota_{\xi^{0,1}}$, so the symbol of $\Box_E$ is $\xi^{0,1} \wedge \iota + \iota\,\xi^{0,1} \wedge = |\xi^{0,1}|^2 = \frac12|\xi|^2$, by the
same computation as in Exercise~\ref{elliptic-operators:exercise-symbol-laplacian}. The rest is the proof of Theorems~\ref{hodge-theorem:theorem-decomposition}
and~\ref{hodge-theorem:theorem-hodge} with $d$, $\delta$ replaced by $\bar\partial_E$, $\bar\partial_E^*$ (Remark~\ref{hodge-theorem:remark-elliptic-complexes}): the only inputs were ellipticity, self-adjointness and
$\bar\partial_E^2 = 0$. Combine with Theorem~\ref{complex-manifolds:theorem-dolbeault}.
\end{proof}

\begin{theorem}[Serre duality]\label{complex-manifolds:theorem-serre-duality}
Let $X$ be a compact complex manifold of dimension $n$ and $E$ a holomorphic vector bundle. The pairing
\[
H^q(X, E) \times H^{n-q}(X, K_X \otimes E^\vee) \to H^n(X, K_X) \xrightarrow{\ \int\ } \CC, \qquad ([\alpha], [\beta]) \mapsto \int_X\alpha \wedge \beta,
\]
computed with Dolbeault representatives of types $(0,q)$ and $(n, n-q)$ and the evaluation $E \otimes E^\vee \to \CC$, is perfect. In particular $h^{p,q} = h^{n-p,n-q}$ for the Hodge numbers.
\end{theorem}

\begin{proof}
The pairing is well defined by Stokes' theorem, since $\int_X\bar\partial(\alpha \wedge \gamma) = \int_Xd(\alpha \wedge \gamma) = 0$ for forms of total type $(n, n-1)$. Choose Hermitian metrics and define the conjugate-linear
operator $\bar*_E \colon \mathcal{A}^{0,q}(E) \to \mathcal{A}^{n,n-q}(E^\vee)$ by $\alpha \wedge \bar*_E\beta = \langle\alpha, \beta\rangle\,\mathrm{vol}$, using the metric to identify $\overline{E}$ with $E^\vee$. As
for the real Hodge star, $\bar\partial^* = -\bar*\bar\partial\bar*$, so $\bar*_E$ maps $\Box_E$-harmonic forms bijectively onto $\Box_{K \otimes E^\vee}$-harmonic forms. For a nonzero harmonic $\alpha$,
$\int_X\alpha \wedge \bar*_E\alpha = \|\alpha\|^2 > 0$, so the pairing is nondegenerate on harmonic representatives, which by Theorem~\ref{complex-manifolds:theorem-dbar-hodge} represent all classes. For the Hodge
numbers take $E = \Omega^p_X$, so $K_X \otimes E^\vee \cong \Omega^{n-p}_X$ by the perfect pairing $\Omega^p \otimes \Omega^{n-p} \to K_X$.
\end{proof}

\section{Examples and the Fr\"olicher inequality}\label{complex-manifolds:section-examples}

\begin{proposition}\label{complex-manifolds:proposition-frolicher-inequality}
For a compact complex manifold, $b_k(X) \leq \sum_{p+q=k}h^{p,q}(X)$, and $\chi(X) = \sum_{p,q}(-1)^{p+q}h^{p,q}(X)$. Equality holds for all $k$ if and only if the Fr\"olicher spectral sequence degenerates at
$E_1$.
\end{proposition}

\begin{proof}
In a spectral sequence of finite-dimensional vector spaces, $\dim E_{r+1}^{p,q} \leq \dim E_r^{p,q}$ and $\sum_{p+q=k}\dim E_\infty^{p,q} = b_k$ (Theorem~\ref{homological-algebra:theorem-filtered-complex}), so
$b_k \leq \sum_{p+q=k}h^{p,q}$, with equality for all $k$ if and only if all differentials $d_r$, $r \geq 1$, vanish. The Euler characteristic is unchanged when passing to cohomology of a differential of
total degree one, so it is the same on every page.
\end{proof}

\begin{example}\label{complex-manifolds:example-hodge-numbers}
\begin{enumerate}
\item For a complex torus $X = \CC^n/\Lambda$, the forms $dz^I \wedge d\bar z^J$ with constant coefficients are harmonic for the flat metric and span each $\mathcal{H}^{p,q}$, so
$h^{p,q} = \binom{n}{p}\binom{n}{q}$, and $\sum_{p+q=k}h^{p,q} = \binom{2n}{k} = b_k$: the Fr\"olicher spectral sequence degenerates.
\item For $\CC\PP^n$, $h^{p,q} = 1$ if $p = q \leq n$ and $0$ otherwise. Since $b_{2k} = 1$ and $b_{2k+1} = 0$, this follows from the Hodge decomposition (Chapter~\ref{kahler}) once one knows
$h^{k,k} \geq 1$, which is witnessed by the class of $\omega_{FS}^k$; see \cite[Chapter 0, Section 7]{griffiths-harris}.
\item For the Hopf surface of Exercise~\ref{complex-manifolds:exercise-hopf-surface}, $b_1 = 1$, $h^{0,1} = 1$ and $h^{1,0} = 0$. So $b_1 = h^{1,0} + h^{0,1}$ (the Fr\"olicher spectral sequence of a compact
complex surface always degenerates), but $h^{1,0} \neq h^{0,1}$: the Hodge symmetry $h^{p,q} = h^{q,p}$ of compact K\"ahler manifolds fails \cite[Chapter IV]{bhpv}.
\end{enumerate}
\end{example}

\begin{example}\label{complex-manifolds:example-riemann-sphere}
The Hodge numbers of $\CC\PP^1$ can be computed directly. First $h^{0,0} = 1$ by Proposition~\ref{complex-manifolds:proposition-no-functions}. A holomorphic $1$-form is $f(z)\,dz$ on the chart $\CC$, with $f$
entire; in the chart $w = 1/z$ at infinity, $dz = -w^{-2}dw$, so $-f(1/w)w^{-2}$ must extend holomorphically to $w = 0$. Hence $|f(z)| \leq C|z|^{-2}$ for $|z| \geq 1$, so $f$ is bounded, hence constant by
Liouville's theorem (Corollary~\ref{complex-analysis:corollary-liouville}), and the constant is $0$. So $h^{1,0} = 0$ by Lemma~\ref{complex-manifolds:lemma-harmonic-closed}. Serre duality
(Theorem~\ref{complex-manifolds:theorem-serre-duality}, with $n = 1$ and $E = \mathcal{O}$ or $E = \Omega^1$) gives $h^{0,1} = h^{1,0} = 0$ and $h^{1,1} = h^{0,0} = 1$. Since $b_0 = b_2 = 1$ and $b_1 = 0$, equality
holds in Proposition~\ref{complex-manifolds:proposition-frolicher-inequality}. Moreover $H^1(\CC\PP^1, \mathcal{O}) = 0$, and $H^2(\CC\PP^1, \mathcal{O}) = 0$ because there are no forms of type $(0,2)$ on a curve; by
Theorem~\ref{complex-manifolds:theorem-exponential}, $c_1 \colon \Pic(\CC\PP^1) \to H^2(\CC\PP^1; \ZZ) \cong \ZZ$ is an isomorphism.
\end{example}

\begin{counterexample}\label{complex-manifolds:counterexample-not-homotopy-invariant}
Dolbeault cohomology is not a homotopy invariant, unlike de Rham cohomology. The plane $\CC$ is contractible, so $H^k_{\mathrm{dR}}(\CC) = 0$ for $k > 0$ and $H^0_{\mathrm{dR}}(\CC) = \CC$. But
$H^{0,0}_{\bar\partial}(\CC)$ and $H^{1,0}_{\bar\partial}(\CC)$ are the spaces of entire functions and of holomorphic $1$-forms $f(z)\,dz$ (there is nothing to divide by in type $(p,0)$), both
infinite-dimensional: they contain the $z^k$ and the $z^k\,dz$. The two computations are reconciled by $\partial \colon H^{0,0}_{\bar\partial}(\CC) \to H^{1,0}_{\bar\partial}(\CC)$, $f \mapsto f'\,dz$, which is
surjective, because every entire function has an entire primitive (Corollary~\ref{complex-analysis:corollary-primitive-convex}), and has kernel the constants. In the Fr\"olicher spectral sequence of Theorem~\ref{complex-manifolds:theorem-holomorphic-de-rham} the first differential is induced
by $\partial$, and here it cancels the two infinite-dimensional groups against each other, leaving $H^0_{\mathrm{dR}}(\CC) = \CC$.
\end{counterexample}

\begin{remark}\label{complex-manifolds:remark-vanishing}
Cohomology of line bundles on projective space is computed by Bott's formula: $H^q(\CC\PP^n, \mathcal{O}(k))$ vanishes for $0 < q < n$ and all $k$, $H^0 = $ homogeneous polynomials of degree $k$, and
$H^n(\CC\PP^n, \mathcal{O}(k)) \cong H^0(\CC\PP^n, \mathcal{O}(-k-n-1))^\vee$ by Serre duality, since $K_{\CC\PP^n} = \mathcal{O}(-n-1)$ (Proposition~\ref{complex-manifolds:proposition-adjunction}).
Vanishing theorems for positive bundles are the subject of Chapter~\ref{kodaira}; the algebraic computation is in Chapter~\ref{schemes}.
\end{remark}

\section{Exercises}\label{complex-manifolds:section-exercises}

\subsection{Complex Manifolds}\label{complex-manifolds:subsection-exercises-complex-manifolds}

\begin{exercise}\label{complex-manifolds:exercise-tori}
Show that every complex torus of dimension $1$ is biholomorphic to $\CC/(\ZZ + \tau\ZZ)$ with $\operatorname{Im}\tau > 0$, and that two such are biholomorphic if and only if the parameters are related by
an element of $\SL_2(\ZZ)$ acting by M\"obius transformations.
\end{exercise}

\begin{exercise}\label{complex-manifolds:exercise-projective-space-charts}
Verify that the transition maps of $\CC\PP^n$ are holomorphic and that $\CC\PP^1$ is biholomorphic to the Riemann sphere $\CC \cup \{\infty\}$, with $S^2$ as underlying smooth manifold.
\end{exercise}

\begin{exercise}\label{complex-manifolds:exercise-hopf-surface}
The \emph{Hopf surface} is $X = (\CC^2 \setminus \{0\})/\langle z \mapsto 2z\rangle$. Show that it is a compact complex manifold diffeomorphic to $S^1 \times S^3$, and deduce that $b_1(X) = 1$; conclude that $X$
carries no K\"ahler metric (Chapter~\ref{kahler}), since odd Betti numbers of compact K\"ahler manifolds are even.
\end{exercise}

\begin{solution}
The action is free and properly discontinuous, so the quotient is a complex manifold, and it is compact because the annulus $1 \leq |z| \leq 2$ maps onto it. The diffeomorphism with $S^1 \times S^3$ comes from
the polar decomposition $\CC^2 \setminus \{0\} \cong S^3 \times \RR_{>0}$, with the action becoming translation on the second factor. Then $b_1(S^1 \times S^3) = 1$ by the K\"unneth formula.
\end{solution}

\begin{exercise}\label{complex-manifolds:exercise-almost-complex-s6}
Show that $S^2$ and $S^6$ admit almost complex structures (for $S^6$ use the octonionic cross product), and that $S^2$'s is integrable.
\end{exercise}

\begin{exercise}\label{complex-manifolds:exercise-residue-infinity}
Compute $\int_{|z| = 2}\frac{z^5}{z^6 - 1}\,dz$ and $\int_{|z| = 2}\frac{dz}{(z - 1)^3(z^2 + 9)}$ (the circle traversed once counterclockwise) by computing a single residue at $\infty$ on $\CC\PP^1$ instead of the residues
inside the circle.
\end{exercise}

\begin{solution}
By Theorem~\ref{complex-manifolds:theorem-global-residue} on $\CC\PP^1$, the integral over the circle, which is $2\pi i$ times the sum of the residues inside it, equals $-2\pi i$ times the sum of the residues
outside it, including $\infty$. For $\omega = z^5\,dz/(z^6 - 1)$ all six poles lie inside. With $w = 1/z$, $\frac{z^5}{z^6 - 1} = \frac{w}{1 - w^6}$ and $dz = -dw/w^2$, so $\omega = -\frac{dw}{w(1 - w^6)}$ has residue $-1$
at $\infty$, and the integral is $2\pi i$. (Directly: each pole $\zeta$ is simple with residue $\zeta^5/(6\zeta^5) = \frac16$.) For the second form the poles $\pm3i$ lie outside the circle and $1$ inside. At
$\infty$ the form is $O(w^3)\,dw$, holomorphic, so the integral is $-2\pi i(\operatorname{Res}_{3i} + \operatorname{Res}_{-3i})$, with
$\operatorname{Res}_{\pm3i} = \frac{1}{(\pm3i - 1)^3(\pm6i)}$. Now $(3i - 1)^3 = 26 - 18i$ and $(3i - 1)^3 \cdot 6i = 108 + 156i$, so $\operatorname{Res}_{3i} = \frac{1}{108 + 156i}$ and
$\operatorname{Res}_{-3i} = \frac{1}{108 - 156i}$ is its conjugate. Their sum is $\frac{216}{108^2 + 156^2} = \frac{216}{36000} = \frac{3}{500}$, and the integral is $-\frac{3\pi i}{250}$. (Directly: the
residue at the triple pole $1$ is $\frac12\frac{d^2}{dz^2}(z^2 + 9)^{-1}$ at $z = 1$, which is $\frac12\bigl(-\frac{2}{100} + \frac{8}{1000}\bigr) = -\frac{3}{500}$.)
\end{solution}

\begin{exercise}\label{complex-manifolds:exercise-abel-torus}
Let $f$ be a nonconstant meromorphic function on $X = \CC/\Lambda$, $\Lambda = \ZZ\omega_1 + \ZZ\omega_2$, viewed as a $\Lambda$-periodic meromorphic function on $\CC$. Let $a_1, \ldots, a_k$ and $b_1, \ldots, b_k$ be
its zeros and poles in a fundamental parallelogram $P$ with vertices $a$, $a + \omega_1$, $a + \omega_1 + \omega_2$, $a + \omega_2$, listed with multiplicity, with $a$ chosen so that no zero or pole lies on
$\partial P$. Show that $\sum_ia_i - \sum_ib_i \in \Lambda$, by computing $\frac{1}{2\pi i}\int_{\partial P}z\frac{f^{\prime}(z)}{f(z)}\,dz$ in two ways. (This is one direction of Abel's theorem for complex tori.)
\end{exercise}

\begin{solution}
At a zero or pole $c$ of order $m$ (negative for a pole), $f^{\prime}/f = m/(z - c) + (\text{holomorphic})$, so $zf^{\prime}/f$ has residue $mc$. The residue theorem in the plane
(Theorem~\ref{complex-analysis:theorem-residues}) gives $\frac{1}{2\pi i}\int_{\partial P}zf^{\prime}/f\,dz = \sum_ia_i - \sum_ib_i$. On the other hand, pair opposite sides: the side from $a + \omega_2$ to
$a + \omega_1 + \omega_2$ is traversed backwards, and on it $f^{\prime}/f$ takes the same values as on the side from $a$ to $a + \omega_1$, so the two sides contribute
\[
\int_a^{a+\omega_1}\bigl(z - (z + \omega_2)\bigr)\frac{f^{\prime}(z)}{f(z)}\,dz = -\omega_2\int_a^{a+\omega_1}\frac{f^{\prime}(z)}{f(z)}\,dz = -2\pi i\,n_1\omega_2,
\]
where $\frac{1}{2\pi i}\int_a^{a+\omega_1}f^{\prime}/f\,dz = n_1$ is the winding number around $0$ of the closed curve $t \mapsto f(a + t\omega_1)$, $t \in [0, 1]$ (closed because $f$ is periodic), an
integer. In the same way the other two sides contribute $2\pi i\,n_2\omega_1$ with $n_2 \in \ZZ$. So $\sum_ia_i - \sum_ib_i = n_2\omega_1 - n_1\omega_2 \in \Lambda$. Applying the same computation to
$f^{\prime}/f$ instead of $zf^{\prime}/f$ gives $k$ zeros and $k$ poles again, in accordance with Corollary~\ref{complex-manifolds:corollary-zeros-poles}.
\end{solution}

\subsection{Holomorphic Vector Bundles}\label{complex-manifolds:subsection-exercises-holomorphic-bundles}

\begin{exercise}\label{complex-manifolds:exercise-sections-ok}
Show that the holomorphic sections of $\mathcal{O}(k)$ on $\CC\PP^n$ are the homogeneous polynomials of degree $k$ in the homogeneous coordinates, and that there are none for $k < 0$.
\end{exercise}

\begin{solution}
A section is given by functions $s_i$ on $U_i$ with $s_i = (x_i/x_j)^{-k}s_j$; the function $F(x) = x_i^ks_i(x/x_i)$ is independent of $i$, holomorphic on $\CC^{n+1} \setminus \{0\}$ and homogeneous of degree
$k$. By Hartogs' theorem (Theorem~\ref{several-complex-variables:theorem-hartogs}) it extends to $\CC^{n+1}$, and a homogeneous holomorphic function of degree $k$ is a polynomial for $k \geq 0$ and zero for
$k < 0$, by looking at its power series.
\end{solution}

\begin{exercise}\label{complex-manifolds:exercise-degree-line-bundle}
Let $X$ be a compact Riemann surface. Show that $\deg L = \int_Xc_1(L)$ defines a homomorphism $\Pic(X) \to \ZZ$, that $\deg\mathcal{O}(D) = \sum_pa_p$ for $D = \sum a_pp$, and that a line bundle of negative
degree has no nonzero holomorphic sections.
\end{exercise}

\begin{exercise}\label{complex-manifolds:exercise-chern-connection-trivial}
Compute the Chern connection and its curvature for the trivial line bundle with metric $e^{-|z|^2}$ on $\CC$, and verify that the curvature is a positive $(1,1)$-form.
\end{exercise}

\begin{exercise}\label{complex-manifolds:exercise-picard-torus}
For a complex torus $X = \CC/\Lambda$ of dimension one, show that $H^1(X, \mathcal{O}_X) \cong \CC$ and deduce that $\Pic^0(X)$ is a complex torus of dimension one.
\end{exercise}

\begin{exercise}\label{complex-manifolds:exercise-zeros-vector-field}
Show that a nonzero holomorphic vector field on $\CC\PP^1$ has exactly two zeros, counted with multiplicity, where the multiplicity of a zero of $f\,\partial_z$ is the order of the zero of $f$. Compare with the Euler
characteristic $\chi(S^2) = 2$ and with $T^{1,0}\CC\PP^1 \cong \mathcal{O}(2)$ (Example~\ref{complex-manifolds:example-tangent-p1}).
\end{exercise}

\begin{solution}
By Example~\ref{complex-manifolds:example-tangent-p1}(1) the field is $f\,\partial_z$ with $f = c_0 + c_1z + c_2z^2 \neq 0$, and near $\infty$ it is $-(c_0w^2 + c_1w + c_2)\,\partial_w$. Let $m = \deg f$. On $\CC$ the zeros of $f$
have total multiplicity $m$, by the fundamental theorem of algebra (Corollary~\ref{complex-analysis:corollary-liouville}(2)) and the factorisation of $f$ into linear factors. At $w = 0$ the coefficient
$c_0w^2 + c_1w + c_2$ has a zero of order $2 - m$, since $c_m$ is its lowest nonzero coefficient, the coefficient of $w^{2-m}$. The total is $m + (2 - m) = 2$. For instance $\partial_z$ has a double zero at $\infty$, and
$z\,\partial_z$ simple zeros at $0$ and $\infty$, the fixed points of the flows $z \mapsto z + t$ and $z \mapsto e^tz$. The number $2$ is the Euler characteristic of $S^2$, as predicted by the Poincar\'e--Hopf theorem
for the real vector field $2\operatorname{Re}(f\,\partial_z)$, and also the degree of $\mathcal{O}(2)$: a nonzero section of $\mathcal{O}(2)$ is a nonzero homogeneous quadric in $x_0, x_1$, which has two zeros on
$\CC\PP^1$ counted with multiplicity.
\end{solution}

\subsection{Dolbeault Cohomology}\label{complex-manifolds:subsection-exercises-dolbeault}

\begin{exercise}\label{complex-manifolds:exercise-riemann-surface}
Let $X$ be a compact Riemann surface of genus $g$. Using Serre duality and the Fr\"olicher inequality, show that $h^{0,0} = h^{1,1} = 1$ and $h^{1,0} = h^{0,1}$, and that $h^{1,0} + h^{0,1} \geq 2g$ with equality
if and only if the Fr\"olicher spectral sequence degenerates (it always does in dimension one, by Chapter~\ref{kahler}).
\end{exercise}

\begin{solution}
$h^{0,0} = 1$ since holomorphic functions are constant (Proposition~\ref{complex-manifolds:proposition-no-functions}); Serre duality gives $h^{1,1} = h^{0,0}$ and $h^{1,0} = h^{0,1}$. The inequality is
Proposition~\ref{complex-manifolds:proposition-frolicher-inequality} for $k = 1$, with $b_1 = 2g$.
\end{solution}

\begin{exercise}\label{complex-manifolds:exercise-dolbeault-plane}
Show that $H^{0,1}_{\bar\partial}(\CC) = 0$, but that $H^{0,1}_{\bar\partial}(\CC^2 \setminus \{0\})$ is infinite-dimensional.
\end{exercise}

\begin{exercise}\label{complex-manifolds:exercise-serre-duality-curve}
For a compact Riemann surface $X$ and a line bundle $L$, show that $H^1(X, L) \cong H^0(X, K_X \otimes L^\vee)^\vee$, and deduce that $H^1(X, L) = 0$ if $\deg L > 2g - 2$.
\end{exercise}

\begin{exercise}\label{complex-manifolds:exercise-euler-characteristic-bundle}
Show that the holomorphic Euler characteristic $\chi(X, E) = \sum_q(-1)^q\dim H^q(X, E)$ of a compact complex manifold is the index of the elliptic operator
$\bar\partial_E + \bar\partial_E^* \colon \mathcal{A}^{0,\mathrm{even}}(E) \to \mathcal{A}^{0,\mathrm{odd}}(E)$.
\end{exercise}
